% Reviewing the history of the internet — ICT Upper Sixth — Cours signé OnBuch+
% Official MINESEC syllabus. Compiler avec : tectonic ICT17-reviewing-the-history-of-the-internet.tex
\newcommand{\DOCMATIERE}{ICT}
\newcommand{\DOCNIVEAU}{Upper Sixth}
\newcommand{\DOCMODULE}{Upper Sixth — ICT}
\newcommand{\DOCLECON}{17}
\newcommand{\DOCTITRE}{Reviewing the history of the internet}
% OnBuch+ — préambule « cours signés » ENRICHI (compilé avec tectonic / XeLaTeX)
% Système visuel « Pop » d'OnBuch+ : fond crème (jamais blanc), bordures encre
% épaisses, ombres « dures » décalées, titres Archivo Black en capitales.
% Version MATHÉMATIQUES : graphes (pgfplots/tikz), boîtes pédagogiques
% (définition, propriété, méthode, attention, le savais-tu, exercice résolu,
% exercice, TP, bilan), page de garde et signature OnBuch+ ancrée en bas.
%
% Macros attendues dans main.tex AVANT \input{preamble} :
%   \DOCMATIERE  \DOCNIVEAU  \DOCMODULE  \DOCLECON (numéro)  \DOCTITRE
\documentclass[11pt]{article}

\usepackage{fontspec}
\usepackage[english]{babel}
\usepackage[a4paper,margin=2cm,top=3.4cm,bottom=2.8cm,headheight=1.4cm,footskip=1.3cm]{geometry}
\usepackage{amsmath,amssymb,amsfonts}
\usepackage{mathtools} % \xmapsto, \coloneqq... souvent utilisés par les modèles
\usepackage{cancel} % simplifications barrées (bilans de forces)
% \abs{z} (module, valeur absolue) et \norm{u} (norme), utilisés par les modèles
\providecommand{\abs}{}\let\abs\relax\DeclarePairedDelimiter{\abs}{\lvert}{\rvert}
\providecommand{\norm}{}\let\norm\relax\DeclarePairedDelimiter{\norm}{\lVert}{\rVert}
\usepackage[table]{xcolor} % [table] : fournit aussi \rowcolors (alternance de lignes)
\usepackage{pagecolor}
\usepackage{tikz}
\usetikzlibrary{shapes.symbols,circuits.logic.US,circuits.logic.IEC,arrows.meta,positioning,calc,shapes.geometric,shapes.misc,shapes.arrows,decorations.pathmorphing,decorations.pathreplacing,decorations.markings,patterns,shadows,mindmap,trees,fit,backgrounds,matrix,intersections,angles,quotes}
\usepackage{pgfplots}
\usepackage[version=4]{mhchem} % équations nucléaires \ce{^{235}_{92}U}
\usepackage[european,siunitx]{circuitikz} % schémas électriques
\pgfplotsset{compat=1.18}
\usepgfplotslibrary{fillbetween,groupplots,polar,statistics,dateplot} % aires entre courbes (calcul intégral)
\usepackage{enumitem}
\usepackage{fancyhdr}
% [most] charge toutes les bibliothèques tcolorbox (listings, minted, raster,
% documentation...) et gonfle inutilement la mémoire TeX sur un document long
% (~300 boîtes) : on ne charge que ce qui est réellement utilisé.
\usepackage[skins,breakable]{tcolorbox}
\usepackage{graphicx}
\usepackage{colortbl}
\usepackage{booktabs}
\usepackage{tabularx}
\usepackage{diagbox} % case d'en-tête diagonale des tableaux à double entrée
% Colonnes extensibles centrée / gauche / droite (C, L, R), souvent utilisées par les modèles
\newcolumntype{C}{>{\centering\arraybackslash}X}
\newcolumntype{L}{>{\raggedright\arraybackslash}X}
\newcolumntype{R}{>{\raggedleft\arraybackslash}X}
\usepackage{array}
\usepackage{multicol}
\usepackage{siunitx}
% parse-numbers=false : accepte aussi \SI{2,5 \times 10^{-3}}{...}
\sisetup{locale=UK,output-decimal-marker={.},group-minimum-digits=4,parse-numbers=false}
\DeclareSIUnit\ohm{\ensuremath{\symup{\Omega}}} % U+2126 absent de la police
\DeclareSIUnit\division{div} % réglages d'oscilloscope (ms/div, V/div)
\DeclareSIUnit\year{yr}\DeclareSIUnit\annum{yr}\DeclareSIUnit\Ma{Ma}\DeclareSIUnit\tr{tr} % tours (vitesse de rotation, ex. \SI{1500}{\tr\per\minute})
\usepackage{qrcode}
% \degree : commande fréquemment utilisée par les modèles pour un angle ou une
% température en dehors de \SI{}{} (non fournie par les paquets chargés).
\providecommand{\degree}{^{\circ}}
% \qed (fin de démonstration), utilisé par les modèles sans amsthm
\providecommand{\qed}{\hfill$\square$}
% \barre : double barre des valeurs interdites dans un tableau de variations (tkz-tab)
\providecommand{\barre}{\ensuremath{\|}}
% environnement proof (amsthm non chargé) utilisé par les modèles
\ifdefined\proof\else\newenvironment{proof}[1][Proof]{\par\noindent\textit{#1.}\ }{\qed\par}\fi
\usepackage{lastpage}
\usepackage{hyperref}
\hypersetup{hidelinks}

% ============================================================
% Palette « Pop » OnBuch+ — mêmes hex que la classe P (plus_ui.dart)
% ============================================================
\definecolor{poporange}{HTML}{F59321}
\definecolor{popdark}{HTML}{E07A0C}
\definecolor{popcream}{HTML}{FFF6E8}
\definecolor{popink}{HTML}{1C1714}
\definecolor{poppurple}{HTML}{7A5AE0}
\definecolor{popgreen}{HTML}{2E9E6A}
\definecolor{poppink}{HTML}{E0457B}
\definecolor{popblue}{HTML}{2F7DE1}
\definecolor{popgold}{HTML}{C98A12}
\definecolor{popmuted}{HTML}{8A7F76}
\definecolor{popline}{HTML}{EADFD2}
\definecolor{popgray}{HTML}{8A7F76}
\colorlet{popgrey}{popgray}
% Alias tolérants (le modèle nomme parfois les couleurs différemment)
\colorlet{popwhite}{white}
\colorlet{popblack}{popink}
\colorlet{popred}{poppink}
\colorlet{popyellow}{popgold}
% Teintes claires pour fonds de tableaux / zones
\colorlet{popblueL}{popblue!10!white}
\colorlet{poporangeL}{poporange!14!white}
\colorlet{popgreenL}{popgreen!12!white}
\colorlet{poppurpleL}{poppurple!12!white}
\colorlet{poppinkL}{poppink!12!white}
\colorlet{popgoldL}{popgold!14!white}

\pagecolor{popcream}

% ============================================================
% Typographie
% ============================================================
\defaultfontfeatures{Ligatures=TeX}
\setmainfont{PlusJakartaSans-Medium.ttf}[Path=../../../fonts/,BoldFont=PlusJakartaSans-Bold.ttf]
% Maths en sans-serif (Fira Math) avec chiffres et lettres droites de Plus
% Jakarta, pour que les formules se fondent dans le texte.
\usepackage[math-style=ISO,bold-style=ISO]{unicode-math}
\setmathfont{FiraMath-Regular.otf}[Path=../../../fonts/]
\setmathfont{PlusJakartaSans-Medium.ttf}[Path=../../../fonts/,range={up/num,up/{Latin,latin},\mathup}]
% (icomma retiré : décimales avec point en anglais)
% Caractères mathématiques Unicode absents des polices de texte (Plus Jakarta,
% Archivo Black) : rendus par leur équivalent mathématique, en texte comme en
% formule — sinon ils s'affichent en carré vide (ex. « Arithmétique dans ℤ »).
\usepackage{newunicodechar}
\newunicodechar{ℝ}{\ensuremath{\mathbb{R}}}
\newunicodechar{ℤ}{\ensuremath{\mathbb{Z}}}
\newunicodechar{ℂ}{\ensuremath{\mathbb{C}}}
\newunicodechar{ℕ}{\ensuremath{\mathbb{N}}}
\newunicodechar{ℚ}{\ensuremath{\mathbb{Q}}}
\newunicodechar{𝔻}{\ensuremath{\mathbb{D}}}
\newunicodechar{α}{\ensuremath{\alpha}}
\newunicodechar{β}{\ensuremath{\beta}}
\newunicodechar{γ}{\ensuremath{\gamma}}
\newunicodechar{δ}{\ensuremath{\delta}}
\newunicodechar{ε}{\ensuremath{\varepsilon}}
\newunicodechar{θ}{\ensuremath{\theta}}
\newunicodechar{λ}{\ensuremath{\lambda}}
\newunicodechar{μ}{\ensuremath{\mu}}
\newunicodechar{σ}{\ensuremath{\sigma}}
\newunicodechar{τ}{\ensuremath{\tau}}
\newunicodechar{φ}{\ensuremath{\varphi}}
\newunicodechar{ψ}{\ensuremath{\psi}}
\newunicodechar{ω}{\ensuremath{\omega}}
\newunicodechar{Σ}{\ensuremath{\Sigma}}
% majuscules produites par \MakeUppercase dans les titres (α-aminés -> Α)
\newunicodechar{Α}{\ensuremath{\alpha}}
\newunicodechar{Β}{\ensuremath{\beta}}
\newunicodechar{Γ}{\ensuremath{\Gamma}}
\newunicodechar{Δ}{\ensuremath{\Delta}}
\newunicodechar{Ε}{\ensuremath{\varepsilon}}
\newunicodechar{Θ}{\ensuremath{\Theta}}
\newunicodechar{Λ}{\ensuremath{\Lambda}}
\newunicodechar{Μ}{\ensuremath{\mu}}
\newunicodechar{Τ}{\ensuremath{\tau}}
\newunicodechar{Φ}{\ensuremath{\Phi}}
\newunicodechar{Ψ}{\ensuremath{\Psi}}
\newunicodechar{Ω}{\ensuremath{\Omega}}
\newunicodechar{↦}{\ensuremath{\mapsto}}
\newunicodechar{∈}{\ensuremath{\in}}
\newunicodechar{∉}{\ensuremath{\notin}}
\newunicodechar{∧}{\ensuremath{\wedge}}
\newunicodechar{⇔}{\ensuremath{\Leftrightarrow}}
\newunicodechar{⟺}{\ensuremath{\Longleftrightarrow}}
\newunicodechar{⇒}{\ensuremath{\Rightarrow}}
\newunicodechar{⟹}{\ensuremath{\Longrightarrow}}
\newunicodechar{∘}{\ensuremath{\circ}}
\newunicodechar{∪}{\ensuremath{\cup}}
\newunicodechar{∩}{\ensuremath{\cap}}
\newunicodechar{≡}{\ensuremath{\equiv}}
\newunicodechar{⊂}{\ensuremath{\subset}}
\newunicodechar{⊥}{\ensuremath{\perp}}
\newunicodechar{∃}{\ensuremath{\exists}}
\newunicodechar{∀}{\ensuremath{\forall}}
\newunicodechar{∛}{\ensuremath{\sqrt[3]{\,}}}
\newunicodechar{∜}{\ensuremath{\sqrt[4]{\,}}}
\newunicodechar{⃗}{\ensuremath{{}^{\to}}}
\newunicodechar{ⁿ}{\ifmmode^{n}\else\textsuperscript{\ensuremath{n}}\fi}
\newunicodechar{ˣ}{\ifmmode^{x}\else\textsuperscript{\ensuremath{x}}\fi}
\newunicodechar{ᵇ}{\ifmmode^{b}\else\textsuperscript{\ensuremath{b}}\fi}
\newunicodechar{ʳ}{\ifmmode^{r}\else\textsuperscript{\ensuremath{r}}\fi}
\newunicodechar{ᵘ}{\ifmmode^{u}\else\textsuperscript{\ensuremath{u}}\fi}
\newunicodechar{ʸ}{\ifmmode^{y}\else\textsuperscript{\ensuremath{y}}\fi}
\newunicodechar{ᵃ}{\ifmmode^{a}\else\textsuperscript{\ensuremath{a}}\fi}
\newunicodechar{ᵉ}{\ifmmode^{e}\else\textsuperscript{\ensuremath{e}}\fi}
\newunicodechar{ᵏ}{\ifmmode^{k}\else\textsuperscript{\ensuremath{k}}\fi}
\newunicodechar{ᵖ}{\ifmmode^{p}\else\textsuperscript{\ensuremath{p}}\fi}
\newunicodechar{ᴺ}{\ifmmode^{N}\else\textsuperscript{\ensuremath{N}}\fi}
\newunicodechar{ᶜ}{\ifmmode^{c}\else\textsuperscript{\ensuremath{c}}\fi}
\newunicodechar{ʰ}{\ifmmode^{h}\else\textsuperscript{\ensuremath{h}}\fi}
\newunicodechar{ᵠ}{\ifmmode^{q}\else\textsuperscript{\ensuremath{q}}\fi}
\newunicodechar{ₙ}{\ifmmode_{n}\else\textsubscript{\ensuremath{n}}\fi}
\newunicodechar{ₐ}{\ifmmode_{a}\else\textsubscript{\ensuremath{a}}\fi}
\newunicodechar{ᵢ}{\ifmmode_{i}\else\textsubscript{\ensuremath{i}}\fi}
\newunicodechar{ᵣ}{\ifmmode_{r}\else\textsubscript{\ensuremath{r}}\fi}
\newunicodechar{ₖ}{\ifmmode_{k}\else\textsubscript{\ensuremath{k}}\fi}
\newunicodechar{ᵦ}{\ifmmode_{\beta}\else\textsubscript{\ensuremath{\beta}}\fi}
\newunicodechar{ₓ}{\ifmmode_{x}\else\textsubscript{\ensuremath{x}}\fi}
\newfontfamily\popdisplay{ArchivoBlack-Regular.ttf}[Path=../../../fonts/]
\newfontfamily\poplogo{SpaceGrotesk-Bold.ttf}[Path=../../../fonts/]
\newfontfamily\popsemi{PlusJakartaSans-SemiBold.ttf}[Path=../../../fonts/]
\newfontfamily\popxbold{PlusJakartaSans-ExtraBold.ttf}[Path=../../../fonts/]
\newfontfamily\popxboldls{PlusJakartaSans-ExtraBold.ttf}[Path=../../../fonts/,LetterSpace=3.0]

\setlist[enumerate]{leftmargin=1.7em,itemsep=5pt,topsep=4pt}
\setlist[itemize]{leftmargin=1.7em,itemsep=4pt,topsep=4pt,label={\color{poporange}\textbullet}}
\setlength{\parindent}{0pt}
\setlength{\parskip}{5pt}
\renewcommand{\arraystretch}{1.25}

% pgfplots : style Pop par défaut
\pgfplotsset{
  every axis/.append style={axis line style={popink,line width=1pt},tick style={popink},
    grid style={popline,line width=0.5pt},label style={font=\small},tick label style={font=\footnotesize}},
  popcurve/.style={poporange,line width=1.8pt,no markers,smooth},
}
% Même style disponible dans un \draw TikZ simple (hors pgfplots)
\tikzset{popcurve/.style={poporange,line width=1.8pt}}
\tikzset{popgrid/.style={popline,very thin}}
\colorlet{popcurve}{poporange} % le nom du style est parfois utilisé comme couleur

% ============================================================
% Pill / squiggle
% ============================================================
\newcommand{\poppill}[3]{%
  \tikz[baseline=-0.55ex]\node[draw=popink,line width=1.2pt,rounded corners=2.6pt,%
    fill=#2,inner xsep=6.5pt,inner ysep=3pt]{%
    \popxboldls\fontsize{7}{7}\selectfont\color{#3}\MakeUppercase{#1}};%
}
\newcommand{\popsquiggle}[1]{%
  \tikz[baseline=-0.4ex]{
    \draw[#1,line width=2.6pt,line cap=round]
      plot [smooth] coordinates {(0,0.09) (0.34,0.01) (0.68,0.21) (1.02,0.01) (1.36,0.10)};
  }%
}
% Mot-clé surligné (marqueur orange sous le texte)
\newcommand{\cle}[1]{{\popxbold\color{popdark}#1}}

% ============================================================
% En-tête / pied
% ============================================================
\pagestyle{fancy}
\fancyhf{}
\renewcommand{\headrulewidth}{0pt}
\renewcommand{\footrulewidth}{0pt}
\lhead{\raisebox{-0.15ex}{\poplogo\fontsize{13}{13}\selectfont\color{popink}OnBuch\color{poporange}+}}
\rhead{\poppill{\DOCMATIERE\ \textbullet\ \DOCNIVEAU\ \textbullet\ Chapter \DOCLECON}{white}{popink}}
\lfoot{\popxbold\fontsize{7.6}{7.6}\selectfont\color{popmuted}\MakeUppercase{Signed course OnBuch+}\ \textbullet\ {\popsemi\DOCTITRE}}
\rfoot{\tikz[baseline=-0.6ex]\node[rounded corners=8.5pt,fill=popink,minimum height=17pt,minimum width=17pt,inner xsep=5pt,inner sep=0]{\popxbold\fontsize{8}{8}\selectfont\color{popcream}\thepage\,/\,\pageref*{LastPage}};}

% ============================================================
% Titres
% ============================================================
\newcommand{\chaptitre}[2]{%
  \begin{tcolorbox}[enhanced,colback=poporange,colframe=popink,boxrule=2.6pt,arc=11pt,
      drop shadow=popink,left=15pt,right=15pt,top=12pt,bottom=11pt,before skip=3pt,after skip=18pt]
    \poppill{#2}{popink}{popcream}\par\vspace{9pt}
    {\raggedright\hyphenpenalty=10000\exhyphenpenalty=10000\popdisplay\fontsize{22}{24}\selectfont\color{popcream}\MakeUppercase{#1}\par}
  \end{tcolorbox}
}
% Section numérotée : pastille orange avec le numéro + titre Archivo
\newcounter{popsec}
\newcommand{\coursec}[1]{%
  \stepcounter{popsec}\setcounter{popsub}{0}%
  \par\vspace{16pt}\noindent
  \begin{minipage}{\linewidth}
  \tikz[baseline=-0.7ex]\node[draw=popink,line width=1.6pt,fill=poporange,rounded corners=4pt,minimum size=22pt,inner sep=0]{\popdisplay\fontsize{12}{12}\selectfont\color{popcream}\thepopsec};%
  \hspace{8pt}{\popdisplay\fontsize{15}{17}\selectfont\color{popink}\MakeUppercase{#1}}\par
  \vspace{4pt}\noindent{\color{poporange}\rule{36pt}{3.6pt}}
  \end{minipage}\par\nopagebreak\vspace{8pt}
}
\newcounter{popsub}
\newcommand{\courssub}[1]{%
  \stepcounter{popsub}%
  \par\vspace{9pt}\noindent{\popxbold\fontsize{11.5}{13}\selectfont\color{popink}{\color{poporange}\thepopsec.\thepopsub}\hspace{6pt}#1}\par\nopagebreak\vspace{4pt}
}

% ============================================================
% Boîtes pédagogiques — même recette : fond blanc, bordure épaisse colorée,
% ombre dure encre, badge plein. Titre optionnel [#1].
% ============================================================
\tcbset{popbase/.style={breakable,enhanced,colback=white,boxrule=2.3pt,arc=9pt,
  shadow={3pt}{-3pt}{0pt}{popink},left=14pt,right=14pt,top=11pt,bottom=11pt,before skip=11pt,after skip=15pt,
  fontupper=\popsemi\fontsize{10.6}{14.5}\selectfont\color{popink},
  fontlower=\popsemi\fontsize{10.6}{14.5}\selectfont\color{popink}}}
% Coche dessinée (le glyphe ✓ n'existe pas dans les polices du document)
\newcommand{\popcheck}[1]{\tikz[baseline=-0.5ex]\draw[#1,line width=1.6pt,line cap=round,line join=round](-0.13,0.0)--(-0.04,-0.09)--(0.13,0.1);}
\providecommand{\coche}{\popcheck{popgreen}}
\providecommand{\resultat}[1]{\par\smallskip\noindent\fcolorbox{popgreen}{popgreenL}{\parbox{\dimexpr\linewidth-2\fboxsep-2\fboxrule\relax}{\textbf{Result.} #1}}\par\smallskip}
\newcommand{\popboxhead}[3]{\poppill{#1}{#2}{white}\ifx\relax#3\relax\else\hspace{6pt}{\popxbold\fontsize{10.6}{12}\selectfont\color{popink}#3}\fi\par\vspace{7pt}}

\newtcolorbox{aretenir}[1][]{popbase,colframe=popgreen,before upper={\popboxhead{Key points}{popgreen}{#1}}}
\newtcolorbox{exemplebox}[1][]{popbase,colframe=poporange,before upper={\popboxhead{Exemple}{poporange}{#1}}}
\newtcolorbox{definition}[1][]{popbase,colframe=popblue,before upper={\popboxhead{Definition}{popblue}{#1}}}
\newtcolorbox{propriete}[1][]{popbase,colframe=poppurple,before upper={\popboxhead{Property}{poppurple}{#1}}}
\newtcolorbox{methode}[1][]{popbase,colframe=popgold,before upper={\popboxhead{Method}{popgold}{#1}}}
\newtcolorbox{attention}[1][]{popbase,colframe=poppink,before upper={\popboxhead{Warning}{poppink}{#1}}}
\newtcolorbox{experience}[1][]{popbase,colframe=popblue,colback=popblueL,before upper={\popboxhead{Experiment}{popblue}{#1}}}
% « Le savais-tu ? » : carte violette pleine (contexte, culture, Cameroun)
\newtcolorbox{savaistu}[1][]{popbase,colback=poppurple,colframe=popink,
  fontupper=\popsemi\fontsize{10.4}{14.2}\selectfont\color{white},
  before upper={\poppill{Did you know?}{white}{poppurple}\ifx\relax#1\relax\else\hspace{6pt}{\popxbold\color{white}#1}\fi\par\vspace{7pt}}}
% Exercice résolu : énoncé en haut, \tcblower puis la solution sur fond crème
\newcounter{popexo}\newcounter{popexr}
\newtcolorbox{exoresolu}[1][]{popbase,colframe=poporange,colbacklower=popcream,
  segmentation style={popink,line width=1.2pt,dashed},
  before upper={\stepcounter{popexr}\popboxhead{Worked example \thepopexr}{poporange}{#1}},
  before lower={\poppill{Solution}{popink}{popcream}\par\vspace{6pt}}}
% Exercice (non résolu en place) — la correction est dans la partie Corrigés
\newtcolorbox{exercice}[1][]{popbase,colframe=popink,
  before upper={\stepcounter{popexo}\popboxhead{Exercise \thepopexo}{popink}{#1}}}
% Corrigé
\newtcolorbox{corrige}[1][]{popbase,colframe=popgreen,colback=popgreenL,
  before upper={\popboxhead{Solution}{popgreen}{#1}}}
% Objectifs / prérequis / bilan
\newtcolorbox{objectifs}{popbase,colframe=popink,colback=poporangeL,
  before upper={\popboxhead{Chapter objectives}{popink}{}}}
\newtcolorbox{prerequis}{popbase,colframe=popink,colback=popblueL,
  before upper={\popboxhead{Prerequisites}{popblue}{}}}
\newtcolorbox{bilan}{popbase,colframe=popink,colback=white,boxrule=2.8pt,
  before upper={\popboxhead{Fiche bilan}{poporange}{L'essentiel à connaître pour l'examen}}}
% Figure encadrée avec légende
\newtcolorbox{popfigure}[1][]{enhanced,breakable=false,colback=white,colframe=popline,boxrule=1.4pt,arc=8pt,
  left=10pt,right=10pt,top=10pt,bottom=8pt,before skip=10pt,after skip=12pt,halign=center,
  lower separated=false,halign lower=center,
  fontlower=\popsemi\fontsize{9}{11.5}\selectfont\color{popmuted},
  #1}
\newcommand{\legende}[1]{\par\vspace{6pt}{\popsemi\fontsize{9}{11.5}\selectfont\color{popmuted}#1}}

% ============================================================
% Signature OnBuch+
% ============================================================
\newcommand{\poplogoinline}[1]{{\poplogo\fontsize{#1}{#1}\selectfont\color{popink}OnBuch\color{poporange}+}}
\newcommand{\signatureonbuch}{%
  \begin{tcolorbox}[enhanced,colback=popink,colframe=popink,boxrule=2.2pt,arc=10pt,
      drop shadow=poporange,left=14pt,right=14pt,top=10pt,bottom=10pt,before skip=0pt,after skip=0pt]
    \tikz[baseline=-0.7ex]\node[circle,fill=popgreen,draw=popcream,line width=1.3pt,minimum size=22pt,inner sep=0]{\popcheck{white}};%
    \hspace{9pt}\begin{minipage}[c]{0.62\linewidth}
      {\popxbold\fontsize{12}{13}\selectfont\color{popcream}Signed\ \textbullet\ {\poplogo OnBuch\color{poporange}+}}\par\vspace{2pt}
      {\raggedright\popsemi\fontsize{8}{10.5}\selectfont\color{popline}\MakeUppercase{OnBuch+ teaching team}\\\MakeUppercase{Follows the official MINESEC syllabus}\par}
    \end{minipage}\hfill
    {\poplogo\fontsize{22}{22}\selectfont\color{popcream}OnBuch\color{poporange}+}
  \end{tcolorbox}%
}

% ============================================================
% Page de garde
% #1 titre · #2 matière · #3 niveau · #4 module · #5 n° leçon · #6 durée ·
% #7 description / objectifs (texte ou itemize)
% ============================================================
\newcommand{\pagegarde}[7]{%
  \thispagestyle{empty}
  \newgeometry{a4paper,margin=1.7cm,top=1.6cm,bottom=1.4cm}%
  \begin{tikzpicture}[remember picture,overlay]
    \fill[poporange!18] ([xshift=-3.2cm,yshift=-3.6cm]current page.north east) circle (4.6cm);
    \fill[poppurple!14] ([xshift=2.4cm,yshift=7.6cm]current page.south west) circle (3.8cm);
  \end{tikzpicture}%
  {\poplogo\fontsize{30}{30}\selectfont\color{popink}OnBuch\color{poporange}+}\hfill
  \raisebox{0.6ex}{\poppill{Signed course \textbullet\ Premium}{popink}{popcream}}\par
  \vspace{1pt}\popsquiggle{poppurple}\par
  \vspace{8pt}
  \begin{tcolorbox}[enhanced,colback=poporange,colframe=popink,boxrule=2.8pt,arc=13pt,
      drop shadow=popink,left=18pt,right=18pt,top=12pt,bottom=13pt,after skip=10pt]
    \poppill{#2 \textbullet\ #3}{popink}{popcream}\hspace{5pt}\poppill{#4}{white}{popink}\par\vspace{8pt}
    {\popdisplay\fontsize{14}{14}\selectfont\color{popink}CHAPTER #5}\par\vspace{4pt}
    {\raggedright\hyphenpenalty=10000\exhyphenpenalty=10000\popdisplay\fontsize{25}{27}\selectfont\color{popcream}\MakeUppercase{#1}\par}
    \vspace{8pt}
    {\popxbold\fontsize{9.5}{9.5}\selectfont\color{popink}\MakeUppercase{Suggested time: #6}}
  \end{tcolorbox}
  \begin{tcolorbox}[enhanced,colback=white,colframe=popink,boxrule=2.2pt,arc=10pt,
      drop shadow=popink,left=16pt,right=16pt,top=10pt,bottom=10pt,after skip=8pt]
    \poppill{In this chapter}{poppurple}{white}\par\vspace{5pt}
    \popsemi\fontsize{9.3}{12.6}\selectfont\color{popink}#7
  \end{tcolorbox}
  \noindent\begin{minipage}[t]{0.487\linewidth}
    \begin{tcolorbox}[enhanced,colback=popgreen,colframe=popink,boxrule=2.2pt,arc=10pt,
        drop shadow=popink,left=11pt,right=11pt,top=10pt,bottom=10pt,height=3.1cm,valign=center]
      \tcbox[colback=white,colframe=popink,boxrule=1.3pt,arc=5pt,boxsep=0pt,nobeforeafter,
        left=3pt,right=3pt,top=3pt,bottom=3pt,valign=center]{\qrcode[height=1.9cm]{https://play.google.com/store/apps/details?id=com.luvvix.onbuch&hl=en}}%
      \hspace{8pt}\begin{minipage}[c]{0.5\linewidth}
        {\popdisplay\fontsize{11}{12.5}\selectfont\color{white}\MakeUppercase{Download OnBuch}}\par\vspace{4pt}
        {\raggedright\popsemi\fontsize{8.4}{11}\selectfont\color{white}Free \textbullet\ Google Play\\AI tutor, past papers, results\par}
      \end{minipage}
    \end{tcolorbox}
  \end{minipage}\hfill
  \begin{minipage}[t]{0.487\linewidth}
    \begin{tcolorbox}[enhanced,colback=poppurple,colframe=popink,boxrule=2.2pt,arc=10pt,
        drop shadow=popink,left=11pt,right=11pt,top=10pt,bottom=10pt,height=3.1cm,valign=center]
      \tikz[baseline=-0.7ex]\node[circle,fill=white,draw=popink,line width=1.3pt,minimum size=1.6cm,inner sep=0]{\popdisplay\fontsize{24}{24}\selectfont\color{poppurple}+};%
      \hspace{8pt}\begin{minipage}[c]{0.6\linewidth}
        {\popdisplay\fontsize{11}{12.5}\selectfont\color{white}\MakeUppercase{Go OnBuch+}}\par\vspace{4pt}
        {\raggedright\popsemi\fontsize{8.4}{11}\selectfont\color{white}All signed courses, unlimited AI tutor, PDF sheets — OnBuch+ tab in the app\par}
      \end{minipage}
    \end{tcolorbox}
  \end{minipage}
  \vspace{8pt}
  \popcontenu
  \vfill
  \signatureonbuch
  \restoregeometry
  \newpage
}

% Rangée « Ce cours contient » (page de garde)
\newcommand{\poptile}[3]{% #1 couleur #2 gros texte #3 légende
  \begin{tcolorbox}[enhanced,colback=#1,colframe=popink,boxrule=2pt,arc=9pt,shadow={2.5pt}{-2.5pt}{0pt}{popink},
      left=6pt,right=6pt,top=9pt,bottom=9pt,halign=center,valign=center,height=1.75cm,width=\linewidth,nobeforeafter]
    {\popdisplay\fontsize{12.5}{13}\selectfont\color{white}\MakeUppercase{#2}}\par\vspace{3pt}
    {\centering\popsemi\fontsize{7.8}{9.5}\selectfont\color{white}#3\par}
  \end{tcolorbox}}
\newcommand{\popcontenu}{%
  \par\noindent{\popxboldls\fontsize{7.5}{7.5}\selectfont\color{popmuted}THIS SIGNED COURSE CONTAINS}\par\vspace{7pt}
  \noindent\begin{minipage}[t]{0.235\linewidth}\poptile{popblue}{Course}{complete, illustrated, step by step}\end{minipage}\hfill
  \begin{minipage}[t]{0.235\linewidth}\poptile{poporange}{Methods}{and worked examples}\end{minipage}\hfill
  \begin{minipage}[t]{0.235\linewidth}\poptile{poppink}{Exercises}{exam style + solutions}\end{minipage}\hfill
  \begin{minipage}[t]{0.235\linewidth}\poptile{popgreen}{Summary sheet}{the essentials to remember}\end{minipage}\par}

% Sommaire
\newtcolorbox{sommaire}{enhanced,colback=white,colframe=popink,boxrule=2.2pt,arc=10pt,
  drop shadow=popink,left=18pt,right=18pt,top=13pt,bottom=13pt,before skip=6pt,after skip=20pt,
  before upper={\poppill{Contents}{poppurple}{white}\par\vspace{9pt}\popsemi\fontsize{10.6}{14.5}\selectfont\color{popink}}}

% Signature de fin de cours — poussée tout en bas de la dernière page
\newcommand{\findecours}{%
  \par\vfill\nopagebreak
  \begin{center}\popsquiggle{poporange}\end{center}\vspace{4pt}
  \signatureonbuch
}
\AtBeginDocument{\def\llbracket{\mathopen{[\mkern-3.5mu[}}\def\rrbracket{\mathclose{]\mkern-3.5mu]}}} % intervalles d'entiers (glyphe ⟦ absent de la police)
% Glyphes absents de Fira Math : redéfinis à partir de glyphes disponibles
\AtBeginDocument{%
  \def\setminus{\mathbin{\backslash}}\let\smallsetminus\setminus
  \def\vdots{\mathord{\vbox{\baselineskip3.2pt\lineskiplimit0pt\kern4pt\hbox{.}\hbox{.}\hbox{.}}}}%
  \def\ddots{\mathinner{\mkern1mu\raise6.5pt\hbox{.}\mkern2mu\raise3.6pt\hbox{.}\mkern2mu\raise0.7pt\hbox{.}\mkern1mu}}%
  \def\triangle{\mathord{\scalebox{0.9}{$\symup{\Delta}$}}}%
  \def\nabla{\mathord{\raisebox{\depth}{\scalebox{1}[-1]{$\symup{\Delta}$}}}}%
  \def\checkmark{\popcheck{popdark}}%
  \def\star{\mathbin{\ast}}\def\bigstar{\mathord{\ast}}%
  \def\diamond{\mathbin{\scalebox{0.6}{$\Diamond$}}}%
  \def\bigcup{\mathop{\mathchoice{\raisebox{-0.3em}{\scalebox{1.7}{$\cup$}}}{\scalebox{1.25}{$\cup$}}{\cup}{\cup}}\displaylimits}%
  \def\bigcap{\mathop{\mathchoice{\raisebox{-0.3em}{\scalebox{1.7}{$\cap$}}}{\scalebox{1.25}{$\cap$}}{\cap}{\cap}}\displaylimits}%
  \def\lesseqqgtr{\mathrel{\lessgtr}}\def\bigoplus{\mathop{\oplus}\displaylimits}%
}
\providecommand{\ssi}{if and only if} % abréviation fréquente
\AtBeginDocument{\providecommand{\wideparen}{\overparen}} % arc de cercle

% Fonctions usuelles que les modèles utilisent sans que amsmath les fournisse
\providecommand{\arsinh}{\operatorname{arsinh}}\providecommand{\arcosh}{\operatorname{arcosh}}
\providecommand{\artanh}{\operatorname{artanh}}\providecommand{\arsech}{\operatorname{arsech}}
\providecommand{\arcsch}{\operatorname{arcsch}}\providecommand{\arcoth}{\operatorname{arcoth}}
\providecommand{\arcsinh}{\operatorname{arsinh}}\providecommand{\arccosh}{\operatorname{arcosh}}
\providecommand{\arctanh}{\operatorname{artanh}}\providecommand{\sech}{\operatorname{sech}}
\providecommand{\csch}{\operatorname{csch}}\providecommand{\arcsec}{\operatorname{arcsec}}
\providecommand{\arccsc}{\operatorname{arccsc}}\providecommand{\arccot}{\operatorname{arccot}}
\providecommand{\sgn}{\operatorname{sgn}}\providecommand{\lcm}{\operatorname{lcm}}
\providecommand{\Var}{\operatorname{Var}}\providecommand{\Cov}{\operatorname{Cov}}
\providecommand{\permil}{\ensuremath{{}^{0}\!/_{00}}}\providecommand{\textperthousand}{\ensuremath{{}^{0}\!/_{00}}}

%% FIGFIX-BEGIN v3 — correctifs globaux des figures (docs/onbuch-generation/tools/figfix)
\usepackage{adjustbox}\usepackage{etoolbox}
% toute figure TikZ/pgfplots plus large que la ligne est réduite à la largeur disponible
\BeforeBeginEnvironment{tikzpicture}{\begin{adjustbox}{max width=\linewidth}}
\AfterEndEnvironment{tikzpicture}{\end{adjustbox}}
% légendes pgfplots lisibles par-dessus les courbes
\pgfplotsset{every axis/.append style={legend style={fill=white,fill opacity=0.92,text opacity=1,draw=black!15,inner sep=3pt}}}
% étiquettes d'arêtes (syntaxe edge["..."]) avec fond blanc
\tikzset{every edge quotes/.append style={fill=white,inner sep=1.2pt}}
% bulles de cartes mentales : texte à la ligne dans la bulle, bulle un peu plus grande
\tikzset{every concept/.append style={text width=2.0cm,align=center,minimum size=2.3cm,inner sep=2pt}}
%% FIGFIX-END
\begin{document}

\pagegarde{Reviewing the history of the internet}{ICT}{Upper Sixth}{Upper Sixth — ICT}{17}{2 periods}{This chapter traces the Internet from its military‑research origins to the global, high‑speed network we use today, and examines the connectivity technologies and technical requirements that make access possible. Mastery of these topics is essential for the GCE Advanced Level ICT examination and for understanding Cameroon’s digital development.
\vspace{4pt}
\begin{multicols}{2}\small
\begin{itemize}
\item Explain the key milestones in the history and development of the Internet (ARPANET, TCP/IP, World Wide Web).
\item Describe the transition from a research network to a commercial, public infrastructure.
\item Identify and compare major Internet connectivity technologies (dial‑up, DSL, cable, fibre, wireless, satellite).
\item Analyse the hardware, software and protocol requirements for establishing an Internet connection.
\item Differentiate IPv4 and IPv6 addressing schemes and their relevance to address exhaustion.
\item Evaluate the impact of bandwidth, latency and reliability on user experience.
\end{itemize}\end{multicols}}

\begin{sommaire}
\begin{multicols}{2}
\begin{enumerate}
\item 1. Origins and Early Development of the Internet
\item 2. Expansion, Commercialisation and the Modern Internet
\item 3. Internet Connectivity Technologies and Media
\item 4. Technical Requirements, Addressing and Security for Internet Access
\item Integration activity
\item Methods and worked examples
\item Exercises
\item Solutions to the exercises
\item Summary sheet
\end{enumerate}
\end{multicols}
\end{sommaire}

\begin{objectifs}
\begin{itemize}
\item Explain the key milestones in the history and development of the Internet (ARPANET, TCP/IP, World Wide Web).
\item Describe the transition from a research network to a commercial, public infrastructure.
\item Identify and compare major Internet connectivity technologies (dial‑up, DSL, cable, fibre, wireless, satellite).
\item Analyse the hardware, software and protocol requirements for establishing an Internet connection.
\item Differentiate IPv4 and IPv6 addressing schemes and their relevance to address exhaustion.
\item Evaluate the impact of bandwidth, latency and reliability on user experience.
\item Apply knowledge to propose suitable connectivity solutions for specific Cameroonian contexts (urban, peri‑urban, rural).
\item Critically assess common misconceptions about Internet speed, bandwidth and security.
\end{itemize}
\end{objectifs}

\begin{prerequis}
\begin{itemize}
\item Basic understanding of computer networks (LAN, WAN, OSI model) from Lower Sixth.
\item Familiarity with binary numbering and data representation.
\item Knowledge of common telecommunications terms (modem, router, ISP).
\item Awareness of Cameroon’s geography and existing telecom operators (MTN, Orange, Camtel).
\item Ability to interpret simple diagrams and tables.
\end{itemize}
\end{prerequis}

\newpage

\coursec{Origins and Early Development of the Internet}

Imagine a student in Yaoundé downloading a 2\,GB video lesson over a 4G connection while a peer in a remote village in the East Region struggles with a 56\,kbps dial‑up line. The stark contrast illustrates how the Internet’s history and the technologies that connect us shape everyday learning opportunities across Cameroon. This section traces the remarkable journey from a Cold‑War research project to the global network that now underpins education, commerce and communication.

\courssub{Cold‑War Context and the Birth of ARPANET}

In the late 1950s, the Soviet launch of Sputnik (1957) spurred the United States to create the Advanced Research Projects Agency (ARPA) in 1958. ARPA’s mandate was to maintain technological superiority, particularly in command‑and‑control systems that could survive a nuclear strike. The prevailing telephone network used \emph{circuit switching} — a dedicated physical path is reserved for the entire call — which is vulnerable because destroying a single switch disconnects all calls passing through it.

\begin{definition}[Packet Switching]
\textbf{Packet switching} breaks data into small, self‑contained units called \cle{packets}. Each packet carries a destination address and is routed independently through the network. Packets may follow different paths and are reassembled at the destination. This provides robustness: if one link fails, packets are rerouted automatically.
\end{definition}

\begin{attention}[Circuit Switching vs Packet Switching]
Do not confuse the two. Circuit switching (traditional telephone) reserves bandwidth end‑to‑end for the session duration; packet switching (Internet) shares links statistically. In exams, you may be asked to compare latency, efficiency and fault tolerance.
\end{attention}

ARPA funded research at UCLA, Stanford Research Institute (SRI), UC Santa Barbara and the University of Utah. On \textbf{29 October 1969}, the first ARPANET message — “LO” (intended “LOGIN”) — was sent from UCLA to SRI. By December 1969, the four‑node network was operational, using the Network Control Protocol (NCP) over 50\,kbps leased lines.

\begin{popfigure}
\begin{tikzpicture}[
    node distance=2.5cm,
    site/.style={circle, draw=popblue, thick, fill=popblueL, minimum width=2.2cm, align=center},
    arrow/.style={->, thick, popdark}
]
\node[site, align=center] (ucla){UCLA\\ (Los Angeles)};
\node[site, right of=ucla, align=center] (sri){SRI\\ (Stanford)};
\node[site, below of=ucla, align=center] (ucsb){UCSB\\ (Santa Barbara)};
\node[site, right of=ucsb, align=center] (utah){Utah\\ (Salt Lake City)};
\draw[arrow] (ucla) -- (sri) node[midway, above] {50 kbps};
\draw[arrow] (ucla) -- (ucsb);
\draw[arrow] (sri) -- (utah);
\draw[arrow] (ucsb) -- (utah);
\end{tikzpicture}
\legende{The four‑node ARPANET in December 1969.}
\end{popfigure}

\courssub{From NCP to TCP/IP — The Protocol Suite that United Networks}

By the early 1970s, multiple packet‑switching networks existed (ARPANET, SATNET, PRNET, CYCLADES in France). They could not interoperate because each used its own host‑to‑host protocol. Vint Cerf and Bob Kahn solved this with the \textbf{Transmission Control Program} (1974, RFC 675), later split into two layers: \textbf{TCP} (reliable end‑to‑end delivery) and \textbf{IP} (addressing and routing).

\begin{propriete}[TCP/IP Layered Architecture (DoD Model)]
The TCP/IP model comprises four layers:
\begin{enumerate}
    \item \textbf{Application Layer} — user‑level protocols (HTTP, FTP, SMTP, DNS).
    \item \textbf{Transport Layer} — TCP (reliable, connection‑oriented) and UDP (unreliable, datagram).
    \item \textbf{Internet Layer} — IP (routing, fragmentation, addressing), ICMP, ARP.
    \item \textbf{Link Layer} — Ethernet, Wi‑Fi, PPP, etc. (physical and data‑link combined).
\end{enumerate}
Each layer adds its own header (encapsulation). This modularity allows innovation at one layer without disturbing others.
\end{propriete}

On \textbf{1 January 1983} (\emph{Flag Day}), ARPANET switched from NCP to TCP/IP, marking the birth of the \textbf{Internet} as a \emph{network of networks}. The decision to make TCP/IP open and royalty‑free was pivotal for global adoption.

\begin{methode}[Tracing Protocol Evolution via RFC Numbers]
\begin{enumerate}
    \item Identify the protocol (e.g., TCP, IP, HTTP).
    \item Visit the RFC Editor website (\texttt{https://www.rfc-editor.org/}).
    \item Search by keyword or number. Key RFCs:
    \begin{itemize}
        \item RFC 791 (IP), RFC 793 (TCP) — 1981 standards.
        \item RFC 1122/1123 — Host/Application requirements.
        \item RFC 2616 (HTTP/1.1), RFC 7540 (HTTP/2).
    \end{itemize}
    \item Note \emph{Obsoletes} and \emph{Updated by} fields to follow the lineage.
    \item For exam purposes, memorise the seminal RFC numbers for TCP, IP, DNS, HTTP.
\end{enumerate}
\end{methode}

\courssub{Naming the Network — The Domain Name System (DNS)}

As the network grew, the central \texttt{HOSTS.TXT} file maintained at SRI became unmanageable. In 1983, Paul Mockapetris designed the \textbf{Domain Name System (DNS)}, a distributed, hierarchical database (RFC 882, RFC 883).

\begin{definition}[Domain Name System (DNS)]
DNS maps human‑readable domain names (e.g., \texttt{www.antic.cm}) to IP addresses (e.g., \texttt{196.201.216.10}). It uses a tree of domains delegated to authoritative name servers. The root zone is served by 13 logical root server clusters (A–M) distributed worldwide.
\end{definition}

The first top‑level domains (TLDs) created in 1984–1985 were:
\begin{center}
\begin{tabularx}{\linewidth}{|l|X|l|}\hline
\textbf{TLD} & \textbf{Purpose} & \textbf{Year} \\ \hline
\texttt{.arpa} & Infrastructure (reverse DNS) & 1984 \\
\texttt{.com} & Commercial organisations & 1985 \\
\texttt{.edu} & Educational institutions (US) & 1985 \\
\texttt{.gov} & US government & 1985 \\
\texttt{.mil} & US military & 1985 \\
\texttt{.org} & Non‑profit organisations & 1985 \\
\texttt{.net} & Network infrastructure & 1985 \\ \hline
\end{tabularx}
\end{center}
Country‑code TLDs (ccTLDs) like \texttt{.cm} (Cameroon, delegated in 1995 and managed by ANTIC) followed ISO 3166‑1 alpha‑2 codes. The \texttt{.int} TLD for international organisations was created later (1988).

\begin{savaistu}[Cameroon's .cm Domain]
The \texttt{.cm} ccTLD is administered by the \textbf{Agence Nationale des Technologies de l'Information et de la Communication (ANTIC)}. It became active in 1995, enabling Cameroonian entities to have a distinct online identity. Today, \texttt{.cm} domains are widely used by government (\texttt{gov.cm}), education (\texttt{edu.cm}), and businesses.
\end{savaistu}

\courssub{The World Wide Web — A Service on Top of the Internet}

By 1989, the Internet supported email, FTP, Usenet, Telnet — but information was scattered and hard to navigate. At CERN, Tim Berners‑Lee proposed a \emph{hypertext} system to link documents across the network. By Christmas 1990 he had implemented:
\begin{itemize}
    \item \textbf{HTML} (HyperText Markup Language) — document structure.
    \item \textbf{HTTP} (HyperText Transfer Protocol) — application‑layer transfer (port 80).
    \item \textbf{URL} (Uniform Resource Locator) — global addressing scheme (\texttt{scheme://host/path}).
    \item The first web browser/editor (\texttt{WorldWideWeb}) and server (\texttt{httpd}).
\end{itemize}

\begin{attention}[Internet vs World Wide Web]
The \textbf{Internet} is the global \emph{infrastructure} (routers, cables, TCP/IP). The \textbf{World Wide Web} is a \emph{service} (HTTP + HTML + URLs) that runs \emph{on top} of the Internet. Confusing them is a classic exam trap. Other Internet services: email (SMTP/IMAP/POP), file transfer (FTP/SFTP), remote login (SSH), VoIP (SIP/RTP).
\end{attention}

The release of \textbf{Mosaic} (1993), the first popular graphical browser, and \textbf{Netscape Navigator} (1994) triggered explosive growth.

\courssub{Expansion, Commercialisation and the Rise of ISPs}

The US National Science Foundation built \textbf{NSFNET} (1986) as a 56\,kbps backbone connecting supercomputer centres. It upgraded to T1 (1.544\,Mbps) in 1988 and T3 (45\,Mbps) in 1991, linking regional networks and universities. Crucially, NSFNET’s \emph{Acceptable Use Policy} initially restricted commercial traffic.

In 1991, the Commercial Internet eXchange (CIX) allowed commercial networks to interconnect. By 1995, NSFNET was decommissioned; the Internet backbone was fully privatized. \textbf{Internet Service Providers (ISPs)} emerged to sell connectivity to businesses and homes — first via dial‑up (modems over PSTN), then ISDN, DSL, cable, fibre and mobile.

In Cameroon, the first ISP licences were granted in the late 1990s (CAMTEL, then private operators). MTN Cameroon and Orange Cameroon later deployed 2G/3G/4G, making mobile the primary Internet access mode for most citizens.

\begin{popfigure}
\begin{tikzpicture}[
    timeline/.style={thick, popblue},
    event/.style={circle, fill=poporange, inner sep=2pt},
    label/.style={text width=3.2cm, align=center, font=\small}
]
% Axis
\draw[timeline] (0,0) -- (15,0);
% Years and events
\foreach \x/\year/\desc in {
    0/1969/ARPANET\\4 nodes,
    2.5/1974/TCP/IP\\design (Cerf\&Kahn),
    5/1983/Flag Day\\TCP/IP standard,
    7/1984/DNS\\first TLDs,
    9/1989/WWW\\proposal (Berners-Lee),
    11/1991/WWW\\public release,
    13/1993/Mosaic\\browser,
    15/1995/NSFNET\\decommissioned
} {
    \draw[event] (\x,0) circle (3pt);
    \node[label, anchor=north] at (\x,-0.5) {\desc};
    \node[label, anchor=south] at (\x,0.5) {\year};
}
\end{tikzpicture}
\legende{Key milestones from ARPANET to the commercial Internet (1969–1995).}
\end{popfigure}

\courssub{Worked Example: Analysing a Historical Scenario}

\begin{exoresolu}[Question]
A 1992 internal memo at a Cameroonian university states: “We will connect to the Internet via NSFNET to access the World Wide Web for research.” Identify three technical or policy inaccuracies in this statement for the year 1992.
\tcblower
\textbf{Solution.}
\begin{enumerate}
    \item \textbf{NSFNET Access Policy:} In 1992, NSFNET’s Acceptable Use Policy still restricted traffic to \emph{research and education}. A university could connect, but commercial use (e.g., consulting for a private firm) was prohibited. The memo implies unrestricted “Internet” access.
    \item \textbf{World Wide Web Availability:} The WWW was publicly released in August 1991, but in 1992 it was still extremely niche — only a handful of servers existed (mostly at CERN and a few universities). Graphical browsers (Mosaic) appeared in 1993. Most Internet traffic in 1992 was email, FTP, Gopher, Telnet. Expecting routine “Web access” is anachronistic.
    \item \textbf{Terminology — Internet vs NSFNET:} NSFNET was a \emph{backbone} funded by NSF; the “Internet” was the collection of interconnected networks (including ARPANET until 1990, CSNET, BITNET gateways, etc.). Saying “connect to the Internet via NSFNET” conflates a backbone with the whole network of networks. A more accurate phrasing: “connect to the Internet through our regional network which peers with NSFNET.”
\end{enumerate}
\end{exoresolu}

\courssub{Key Points to Retain}

\begin{aretenir}[Summary of Section 1]
\begin{itemize}
    \item \textbf{1969:} ARPANET (4 nodes) demonstrates packet switching — robust, decentralized.
    \item \textbf{1974:} Cerf \& Kahn publish TCP/IP design (RFC 675, later RFC 791/793).
    \item \textbf{1 Jan 1983:} Flag Day — ARPANET adopts TCP/IP; the Internet is born.
    \item \textbf{1984–85:} DNS introduced; first TLDs (\texttt{.arpa}, \texttt{.com}, \texttt{.edu}, \texttt{.gov}, \texttt{.mil}, \texttt{.org}, \texttt{.net}).
    \item \textbf{1989–1991:} Tim Berners‑Lee invents the WWW (HTML, HTTP, URL) at CERN.
    \item \textbf{1993:} Mosaic browser popularises the Web.
    \item \textbf{1995:} NSFNET decommissioned; Internet backbone fully commercial; ISPs proliferate.
    \item \textbf{Distinction:} Internet = infrastructure (TCP/IP); Web = application service (HTTP/HTML).
\end{itemize}
\end{aretenir}

\begin{exercice}[Exercise 1]
\begin{enumerate}
    \item Explain why packet switching is more resilient than circuit switching for a military command network.
    \item List the four layers of the TCP/IP model and give one protocol example per layer.
    \item What problem did DNS solve, and why was a hierarchical design chosen?
    \item Distinguish between the Internet and the World Wide Web. Give two examples of Internet services that are \emph{not} the Web.
    \item In what year did the NSFNET backbone cease operation, and what did this signify for Internet governance?
\end{enumerate}
\end{exercice}

\begin{corrige}[Exercise 1]
\begin{enumerate}
    \item Packet switching does not reserve a dedicated path; packets are routed independently. If a switch or link is destroyed, subsequent packets find alternative routes. Circuit switching fails completely if any link in the dedicated circuit is broken.
    \item Application (HTTP), Transport (TCP), Internet (IP), Link (Ethernet).
    \item DNS replaced the centrally maintained \texttt{HOSTS.TXT} which did not scale. Hierarchy allows delegation: each zone manages its own subdomains, enabling distributed administration and caching.
    \item Internet = global network of networks running TCP/IP. Web = information system using HTTP/HTML/URLs over the Internet. Other services: Email (SMTP), File Transfer (FTP), Remote Login (SSH).
    \item 1995. It marked the transition from a US government‑funded research backbone to a fully privatized, commercially operated Internet infrastructure.
\end{enumerate}
\end{corrige}

\coursec{The Web Revolution and the Modern Internet}

In the previous section we saw how the \cle{internet} grew from the military-academic network ARPANET (1969) into a worldwide infrastructure unified by the TCP/IP protocol suite (1983). By the end of the 1980s, the internet existed and worked well — but it was used almost exclusively by scientists and engineers typing cryptic commands on text terminals. Sending a file required FTP commands; reading news required Usenet; there was no simple way to jump from one document to another. The revolution that transformed this expert tool into the everyday medium we know today was not a new network: it was a new \emph{service} built on top of the existing network. That service is the \cle{World Wide Web}, and understanding the difference between the two is one of the most frequently tested points at the GCE Advanced Level.

\courssub{The Internet Is Not the Web}

Imagine the road network of Cameroon: the tarred roads from Douala to Yaound\'e and Bamenda existed before any particular business was built along them. The roads are the \emph{infrastructure}; the banks, restaurants and delivery services built beside them are \emph{services} that use the roads. Closing one restaurant does not destroy the road, and the road can carry many different kinds of traffic.

The same logic applies here. The \cle{internet} is the global infrastructure: millions of interconnected networks, routers, cables and protocols (TCP/IP) that transport packets of data. The \cle{World Wide Web} is only one service among several that travel over this infrastructure, exactly as e-mail, file transfer (FTP) or video streaming do.

\begin{definition}[The World Wide Web]
The \cle{World Wide Web} (WWW or simply ``the Web'') is a system of interlinked \cle{hypertext} documents and resources, identified by addresses called URLs and accessed via the internet using application programs called \cle{browsers}. It was invented by the British computer scientist \cle{Tim Berners-Lee} at CERN in 1989. The Web is a \emph{service} that runs on the internet; it is not the internet itself.
\end{definition}

\begin{popfigure}
\begin{center}
\begin{tikzpicture}[font=\small]
% infrastructure layer
\draw[fill=popblueL, draw=popblue, rounded corners=2pt] (0,0) rectangle (10,1.6);
\node[align=center] at (5,0.8) {\textbf{INTERNET (infrastructure)}\\ TCP/IP, routers, fibre-optic and submarine cables, satellites};
% services layer
\draw[fill=popgreenL, draw=popgreen, rounded corners=2pt] (0,2.1) rectangle (2.3,3.3);
\node[align=center] at (1.15,2.7) {\textbf{Web}\\ HTTP};
\draw[fill=popgreenL, draw=popgreen, rounded corners=2pt] (2.55,2.1) rectangle (4.85,3.3);
\node[align=center] at (3.7,2.7) {\textbf{E-mail}\\ SMTP};
\draw[fill=popgreenL, draw=popgreen, rounded corners=2pt] (5.1,2.1) rectangle (7.4,3.3);
\node[align=center] at (6.25,2.7) {\textbf{FTP}\\ files};
\draw[fill=popgreenL, draw=popgreen, rounded corners=2pt] (7.65,2.1) rectangle (10,3.3);
\node[align=center] at (8.82,2.7) {\textbf{Streaming}\\ video};
% arrows
\foreach \x in {1.15,3.7,6.25,8.82}{\draw[->, thick, popmuted] (\x,2.1) -- (\x,1.65);}
% users
\draw[fill=poporange!25, draw=poporange, rounded corners=2pt] (0,3.8) rectangle (10,4.6);
\node at (5,4.2) {\textbf{USERS} with browsers, mail clients, apps};
\foreach \x in {1.15,3.7,6.25,8.82}{\draw[->, thick, popmuted] (\x,3.8) -- (\x,3.35);}
\end{tikzpicture}
\end{center}
\legende{Layered view: the internet is the transport infrastructure; the Web, e-mail, FTP and streaming are services running on top of it.}
\end{popfigure}

\begin{attention}[The classic exam trap]
A very common mistake in GCE scripts is to write: \emph{``The internet was invented in 1989 by Tim Berners-Lee.''} This is \textbf{wrong}. What Tim Berners-Lee invented in 1989 is the \textbf{World Wide Web}. The \textbf{internet} is much older: ARPANET dates from \textbf{1969}, and the modern internet based on TCP/IP dates from \textbf{1983}. Always write ``the Web'' when you mean 1989, and ``the internet'' when you mean 1969/1983. Confusing the two can cost you the full mark of a definition question.
\end{attention}

\courssub{1989--1991: Tim Berners-Lee Invents the Web at CERN}

CERN, the European Organisation for Nuclear Research near Geneva, hosted thousands of physicists from different countries, each with different computers and incompatible document systems. Tim Berners-Lee observed that scientists wasted enormous time simply \emph{finding} information scattered across machines. His idea, proposed in March 1989, was elegant: let any document contain \cle{hyperlinks} — clickable references that jump directly to another document, wherever it is stored in the world.

To make this work, he created three technologies that remain the foundation of the Web today:

\begin{itemize}
\item \cle{HTML} (HyperText Markup Language): the language used to write web pages, marking titles, paragraphs, images and, above all, hyperlinks.
\item \cle{HTTP} (HyperText Transfer Protocol): the set of rules by which a browser \emph{requests} a page from a web server and the server \emph{delivers} it.
\item \cle{URL} (Uniform Resource Locator): the unique address of every resource on the Web, for example an address beginning with \emph{http://} followed by a server name and a file path.
\end{itemize}

By 1990--1991 Berners-Lee had also written the first web server and the first browser, which he named \emph{WorldWideWeb} — a program that could both \emph{read} and \emph{edit} pages. The first website, hosted at CERN, simply explained what the Web was. Crucially, CERN decided in 1993 to make the Web technology freely available to everyone, without royalties. This open decision is a major reason the Web spread so fast.

\begin{savaistu}[Key persons to remember for the exam]
\begin{itemize}
\item \cle{Tim Berners-Lee}: invented the World Wide Web (HTML, HTTP, URL) at CERN in 1989.
\item \cle{Ray Tomlinson}: sent the first networked \cle{e-mail} in 1971 and chose the \emph{@} symbol to separate the user name from the machine name.
\item \cle{Marc Andreessen}: co-author of the \cle{Mosaic} browser (1993) and co-founder of \cle{Netscape} (1994), which brought the Web to the general public.
\end{itemize}
\end{savaistu}

\courssub{1993--1998: Browsers for Everyone, Commercialisation and Search Engines}

The early Web was still text-based and known only in universities. Two events opened it to the public. In \textbf{1993}, Marc Andreessen and Eric Bina at the University of Illinois released \cle{Mosaic}, the first popular \emph{graphical} browser: it displayed images inside pages and could be driven with a mouse, making the Web attractive to non-specialists. In \textbf{1994} Andreessen co-founded \cle{Netscape}, whose Navigator browser dominated the market and triggered an explosion of websites. Microsoft answered with Internet Explorer (1995), launching the ``browser wars''.

Businesses quickly understood that a website was a shop window open to the whole planet. The late 1990s saw the \cle{dot-com boom}: investors poured money into any company whose name ended in ``.com'', stock values soared, and thousands of online businesses were created. Many had no realistic revenue, and the bubble burst around 2000--2001 — but the commercial Web survived and matured (Amazon and Google, founded in the 1990s, are famous survivors).

As the number of pages grew into the millions, finding information became a problem, solved by \cle{search engines}: \emph{Yahoo!} (1994, initially a human-edited directory), \emph{AltaVista}, and then \cle{Google} (1998), whose PageRank algorithm ranked pages by the links pointing to them and quickly became the world reference.

\courssub{The 2000s: Web 2.0 — the User Becomes a Producer}

The first Web was largely \emph{read-only}: a few organisations published pages and everyone else consumed them. From the mid-2000s, a new generation of sites inverted this relationship. Ordinary users could publish, comment, share and collaborate without knowing HTML. This transformation is called \cle{Web 2.0}.

\begin{definition}[Web 2.0]
\cle{Web 2.0} designates the \emph{participatory} or \emph{social} Web, in which users are no longer passive readers but active \cle{content producers}: they publish articles and photos, edit shared documents, comment and interact. Typical Web 2.0 services include \cle{blogs}, \cle{wikis} (e.g.\ Wikipedia), \cle{social networks} (Facebook, launched 2004), video-sharing platforms (YouTube, 2005) and messaging applications (WhatsApp, 2009).
\end{definition}

\begin{exemplebox}[Web 1.0 versus Web 2.0 in daily life]
Reading an online newspaper article in 1999 was Web 1.0: you could only read. Today a student in Buea can write a blog post about local agriculture, upload a tutorial video to YouTube, contribute to the Wikipedia article on Mount Cameroon, and discuss it all in a WhatsApp group — that is Web 2.0. The same infrastructure (the internet) carries both; what changed is \emph{who produces the content}.
\end{exemplebox}

\courssub{The Mobile Internet: the Web in Your Pocket}

The next revolution was mobility. Early mobile phones offered only voice and SMS. The arrival of \cle{3G} networks (2000s) and then \cle{4G} (2010s) provided data rates high enough for comfortable browsing, e-mail and video on the move. The launch of Apple's iPhone (2007) and of Android phones popularised the \cle{smartphone}: a pocket computer with a touch screen, a real browser and \cle{apps} — small dedicated applications for banking, maps, messaging or social media. The internet ceased to be a place you ``went to'' at a cybercaf\'e; it became something permanently in your pocket. In Cameroon, as in most of Africa, the majority of people experience the internet \emph{primarily} through a mobile phone, often using mobile money services such as MTN Mobile Money or Orange Money that rely entirely on this connectivity.

\begin{popfigure}
\begin{center}
\begin{tikzpicture}[font=\scriptsize]
\draw[->, thick, popdark] (0,0) -- (12.6,0);
\foreach \x/\y in {0/1989, 2.4/1993, 4.2/1998, 6.6/2004, 8.4/2007, 10.6/2010}{
  \draw[popdark] (\x,0.08) -- (\x,-0.08) node[below] {\y};
}
% world events above
\node[above, align=center, text width=1.7cm, popblue] at (0,0.15) {Berners-Lee invents the Web};
\node[above, align=center, text width=1.7cm, popblue] at (2.4,0.15) {Mosaic; Netscape 1994};
\node[above, align=center, text width=1.7cm, popblue] at (4.2,0.15) {Google; dot-com boom};
\node[above, align=center, text width=1.7cm, popblue] at (6.6,0.15) {Web 2.0: Facebook, YouTube};
\node[above, align=center, text width=1.7cm, popblue] at (8.4,0.15) {iPhone: mobile internet};
\node[above, align=center, text width=1.7cm, popblue] at (10.6,0.15) {Apps and 3G spread};
% Cameroon events below
\node[below, align=center, text width=1.9cm, poporange] at (3.3,-0.45) {1990s: first internet links in Cameroon};
\node[below, align=center, text width=1.9cm, poporange] at (5.6,-0.45) {2002: SAT-3 cable lands};
\node[below, align=center, text width=1.9cm, poporange] at (8.0,-0.45) {MTN, Orange, Nexttel mobile data};
\node[below, align=center, text width=1.9cm, poporange] at (10.9,-0.45) {2012: WACS and ACE cables};
\draw[poporange, dashed] (3.3,-0.35) -- (3.3,-0.1);
\draw[poporange, dashed] (5.6,-0.35) -- (5.6,-0.1);
\draw[poporange, dashed] (8.0,-0.35) -- (8.0,-0.1);
\draw[poporange, dashed] (10.9,-0.35) -- (10.9,-0.1);
\end{tikzpicture}
\end{center}
\legende{Timeline 1989--2012: the Web revolution worldwide (blue, above) and milestones of the internet in Cameroon (orange, below).}
\end{popfigure}

\courssub{A Brief History of the Internet in Cameroon}

Cameroon's first internet connections appeared in the \textbf{1990s}, initially through slow and expensive satellite links and dial-up telephone lines, mostly serving universities, NGOs and a few companies in Douala and Yaound\'e. The state operator \cle{CAMTEL} (Cameroon Telecommunications), created in 1998, inherited the national telephone network and progressively built the national fibre-optic \cle{backbone} that carries traffic between cities and to the coast.

A decisive step was the landing of \emph{submarine fibre-optic cables}, which connect Cameroon to the global internet at high capacity and lower cost than satellites:

\begin{itemize}
\item \cle{SAT-3/WASC/SAFE} (around 2002): the first submarine cable serving Cameroon, landing at Douala.
\item \cle{WACS} (West Africa Cable System, 2012) and \cle{ACE} (Africa Coast to Europe, 2012): newer, higher-capacity cables that multiplied available bandwidth and helped reduce prices.
\end{itemize}

On the retail side, mobile operators \cle{MTN} and \cle{Orange} (from 1999--2000), later joined by \cle{Nexttel} (2014), democratised access: first through 2G/EDGE, then 3G and 4G, making mobile internet the dominant mode of access for Cameroonians. The regulator (ART) and the security agency \cle{ANTIC} supervise the sector, while e-government services (online tax declarations, exam results) progressively move public services onto the Web.

\courssub{Current Trends: Internet of Things and Web 3.0 (Exam Extra)}

\begin{savaistu}[Beyond the syllabus — short awareness]
Two trends are frequently mentioned in exam ``open'' questions:
\begin{itemize}
\item The \cle{Internet of Things} (IoT): everyday objects — meters, vehicles, farm sensors, televisions — are connected to the internet and exchange data without human intervention.
\item \cle{Web 3.0}: an evolving term for a more intelligent and decentralised Web, combining the \emph{semantic web} (data understandable by machines), artificial intelligence and, in some visions, blockchain technologies.
\end{itemize}
A single well-phrased sentence on each is enough to earn the bonus mark; do not develop them at length.
\end{savaistu}

\courssub{Worked Example and Examination Method}

\begin{exoresolu}[Distinguishing internet and Web in a GCE question]
\textbf{Statement.} (a) Define the World Wide Web. (b) A candidate writes: ``The internet was created in 1989 at CERN.'' Identify and correct the two errors in this statement.
\tcblower
\textbf{Solution.} (a) The World Wide Web is a system of interlinked hypertext documents and resources, identified by URLs and accessed through the internet using browsers; it was invented by Tim Berners-Lee at CERN in 1989. (b) \emph{First error}: what was created in 1989 at CERN is the \textbf{Web}, a service, not the internet. \emph{Second error}: the internet is far older — it descends from \textbf{ARPANET (1969)} and took its modern form with the adoption of \textbf{TCP/IP in 1983}. Corrected sentence: ``The \emph{World Wide Web} was created in 1989 at CERN by Tim Berners-Lee; it runs on the internet, whose origins go back to ARPANET in 1969.''
\end{exoresolu}

\begin{methode}[Structuring a ``history of the internet'' essay for the GCE]
When a question asks you to trace the history of the internet, organise your answer in three clearly labelled eras:
\begin{enumerate}
\item \textbf{ARPANET era (1969--1982):} military and academic packet-switching network in the USA; first e-mail (Tomlinson, 1971).
\item \textbf{TCP/IP and NSFNET era (1983--1990):} adoption of the TCP/IP protocols (1983), the internet becomes a network of networks open to universities.
\item \textbf{Web era (1989 to today):} Berners-Lee invents the Web (1989); Mosaic/Netscape open it to the public (1993--94); commercialisation and search engines (late 1990s); Web 2.0 and mobile internet (2000s); in Cameroon, CAMTEL backbone, SAT-3/WACS/ACE cables and mobile operators.
\end{enumerate}
Finish with one sentence on current trends (IoT, Web 3.0). This chronological skeleton guarantees full coverage and impresses the examiner.
\end{methode}

\begin{aretenir}[Key points of this section]
\begin{itemize}
\item The \cle{internet} is the infrastructure (1969/1983); the \cle{Web} is a service on top of it, invented by \cle{Tim Berners-Lee} at CERN in \cle{1989} using HTML, HTTP and URLs.
\item \cle{Mosaic} (1993) and \cle{Netscape} (1994, Marc Andreessen) made the Web graphical and public; the late 1990s brought commercialisation, the dot-com boom and search engines (Google, 1998).
\item \cle{Web 2.0} (2000s) turned users into content producers (blogs, wikis, Facebook, YouTube, WhatsApp); smartphones, 3G/4G and apps made the internet mobile.
\item In Cameroon: first links in the 1990s, CAMTEL backbone, submarine cables SAT-3 (2002), WACS and ACE (2012), and mobile internet through MTN, Orange and Nexttel.
\item Current trends to mention briefly: the Internet of Things and Web 3.0.
\end{itemize}
\end{aretenir}

\begin{exercice}[Check your understanding]
\begin{enumerate}
\item Give two services, other than the Web, that run on the internet, and name the protocol associated with each.
\item Explain, in three sentences, why the statement ``Google invented the internet'' contains two mistakes.
\item Construct a table with two columns (Web 1.0, Web 2.0) and three rows (role of the user, example services, approximate period).
\item Describe the role of submarine cables in the development of the internet in Cameroon, naming at least two of them.
\end{enumerate}
\end{exercice}

\coursec{Internet Connectivity Technologies}

In the previous section we saw how the World Wide Web transformed the internet from a tool for researchers into a global public space. But a web browser is useless if the computer cannot actually \emph{reach} the network. Between your laptop in Bamenda and a web server in London, there must be a physical path for data: a telephone line, a fibre optic cable, a radio wave to a satellite, or a signal from a mobile phone mast. The technologies that provide this path are called \cle{internet connectivity technologies}. Choosing between them is a daily practical question in Cameroon: should a cybercafé in Ngaoundéré invest in fibre, in a 4G router, or in a satellite dish? This section gives you the knowledge to answer such questions rigorously.

\courssub{Classification of Connection Types}

Internet connection technologies are usually classified into four broad families:

\begin{itemize}
\item \cle{Dial-up} connections: narrowband, established on demand over the ordinary telephone line.
\item \cle{Broadband wired} connections: high-speed, permanent connections over telephone lines (ADSL), TV cables, or fibre optics.
\item \cle{Wireless} connections: local wireless access (Wi-Fi hotspots, fixed wireless) and satellite links.
\item \cle{Mobile} connections: internet through the cellular phone network (2G, 3G, 4G, 5G).
\end{itemize}

\begin{definition}[Broadband]
\cle{Broadband} is a high-speed, \emph{always-on} internet connection, as opposed to dial-up, which is slow and must be established (``dialled'') each time it is used. With broadband, the connection is permanent: the moment you open your browser, you are online.
\end{definition}

\begin{definition}[Bandwidth]
The \cle{bandwidth} of a connection is the maximum rate at which data can be transferred through it. It is measured in bits per second (bps) and its multiples: kilobits per second (kbps), megabits per second (Mbps) and gigabits per second (Gbps). A higher bandwidth means web pages load faster, videos play smoothly, and files download quickly.
\end{definition}

\begin{definition}[Latency]
The \cle{latency} of a connection is the time taken for a small piece of data to travel from the sender to the receiver and back (round-trip time), usually measured in milliseconds (ms). Low latency is essential for real-time applications such as video calls and online games; high latency is tolerable for e-mail and ordinary browsing.
\end{definition}

\begin{attention}[Bandwidth is not ``speed'' in the everyday sense]
Students often say a connection is ``fast'' when they mean it has high bandwidth. Strictly, bandwidth measures \emph{how much} data passes per second, not how quickly a single bit travels. Also note the units: internet speeds are given in \textbf{bits} per second (Mbps), while file sizes are in \textbf{Bytes} (1 Byte = 8 bits). A 10 MB file on an 8 Mbps line takes about 10 seconds ($10 \times 8 = 80$ megabits, divided by 8 megabits per second), not 1.25 seconds.
\end{attention}

\courssub{Dial-up: the Historical Connection}

The earliest home internet connections used the ordinary telephone network. A \cle{dial-up modem} (the word comes from MOdulator--DEModulator) converts the computer's digital signals into audio tones that travel over the telephone line, and back again at the other end. To connect, the modem literally \emph{telephoned} the Internet Service Provider (ISP), which is why old modems made characteristic beeping and hissing sounds.

The fastest dial-up modems reached only \textbf{56 kbps} (0.056 Mbps). Two serious limitations followed:

\begin{itemize}
\item The telephone line was \emph{occupied}: nobody could make or receive a phone call while the computer was online.
\item Speeds were so low that loading a single photograph could take a minute; video was impossible.
\end{itemize}

Despite this, dial-up has enormous historical importance: it carried the internet into homes and schools during the 1990s, and in Cameroon the first cybercafés of the early 2000s often ran on dial-up or similar low-speed links.

\courssub{ADSL: Broadband over the Telephone Line}

\cle{ADSL} (Asymmetric Digital Subscriber Line) is a broadband technology that reuses the \emph{same} copper telephone line as dial-up, but in a clever way: it transmits data on high frequencies that the human voice does not use. A small device called a \cle{splitter} (or filter) separates the voice frequencies from the data frequencies, so the telephone and the internet can be used \emph{at the same time}.

The word \emph{asymmetric} means the download speed (from the internet to you) is much higher than the upload speed (from you to the internet). Typical ADSL offers 8 to 24 Mbps download and about 1 Mbps upload. This suits most home users, who download far more than they upload.

\begin{attention}[Distance matters with ADSL]
ADSL speed decreases strongly with the length of the copper line between the subscriber and the telephone exchange. A customer 5 km from the exchange may get only a fraction of the advertised speed. In exams, do not quote the ``advertised'' speed as guaranteed.
\end{attention}

\courssub{Cable and Fibre Optics (FTTH)}

\textbf{Cable internet} delivers broadband through the same coaxial cable used for cable television, typically reaching tens to a few hundred Mbps. It is common in some countries but little deployed in Cameroon.

\textbf{Fibre optics} goes much further. In a fibre optic cable, data is encoded as \emph{pulses of light} travelling inside a very thin glass fibre, kept inside by total internal reflection. When fibre reaches the subscriber's home, the service is called \cle{FTTH} (Fibre To The Home).

\begin{propriete}[Advantages of fibre optics]
Fibre optics offers the \emph{highest speeds} (commonly 100 Mbps to 1 Gbps and beyond) and the \emph{lowest interference} of all wired technologies, because data travels as light pulses in glass fibres: light is immune to the electromagnetic interference that disturbs electrical signals in copper wires, and signal loss over distance is very small.
\end{propriete}

In Cameroon, fibre is deployed mainly in Yaoundé and Douala, where operators such as CAMTEL and MTN offer FTTH to homes and businesses. The national fibre backbone also links major towns, but extending fibre to every quarter remains costly.

\courssub{Satellite Internet}

Where no cable or telephone line reaches --- a village in the Far North, a forest camp in the East --- \cle{satellite internet} may be the only option. The principle is simple to picture:

\begin{enumerate}
\item The user installs a small \emph{dish antenna} pointed at a \emph{geostationary satellite} fixed high above the equator.
\item The dish sends the user's requests up to the satellite, which relays them down to the ISP's \emph{ground station} connected to the internet.
\item Replies travel back along the reverse path: ground station $\rightarrow$ satellite $\rightarrow$ dish $\rightarrow$ computer.
\end{enumerate}

\begin{popfigure}
\begin{center}
\begin{tikzpicture}[scale=0.95, every node/.style={font=\small}]
% Ground
\fill[popgreenL] (-0.5,0) rectangle (12.5,-0.3);
\draw[popline] (-0.5,0) -- (12.5,0);
% User dish
\draw[fill=popblueL, draw=popblue] (1.6,0) -- (2.4,0) -- (2.0,0.5) -- cycle;
\draw[popblue, thick] (2.0,0.5) -- (2.35,1.15);
\draw[popblue, thick, fill=popblueL] (2.05,1.0) arc (200:340:0.35);
\node[align=center] at (2.0,-0.7) {User dish\\(home / school)};
% Satellite
\draw[fill=popgold, draw=popdark] (5.6,5.2) rectangle (6.6,5.8);
\draw[popdark] (5.0,5.35) -- (5.6,5.5); \draw[popdark] (6.6,5.5) -- (7.2,5.35);
\draw[fill=popblueL, draw=popblue] (4.4,5.2) rectangle (5.0,5.65);
\draw[fill=popblueL, draw=popblue] (7.2,5.2) rectangle (7.8,5.65);
\node[align=center] at (6.1,6.3) {Geostationary satellite};
% ISP ground station
\draw[fill=popblueL, draw=popblue] (9.8,0) rectangle (11.0,1.0);
\draw[popblue, thick] (10.4,1.0) -- (10.4,1.5);
\draw[popblue, thick, fill=popblueL] (10.4,1.5) arc (180:360:0.35);
\node[align=center] at (10.4,-0.7) {ISP ground station};
% Internet cloud (drawn as a labelled ellipse)
\draw[fill=popblueL, draw=popink, thick] (10.4,3.4) ellipse (1.3 and 0.65);
\node at (10.4,3.4) {\textbf{Internet}};
% Links
\draw[->, very thick, poporange] (2.5,1.35) -- (5.5,5.35) node[midway, above left, popdark]{uplink};
\draw[->, very thick, poporange] (6.7,5.35) -- (10.15,1.9) node[midway, above right, popdark]{downlink};
\draw[<->, thick, popink] (10.4,2.2) -- (10.4,2.7);
\end{tikzpicture}
\end{center}
\legende{Schematic of a satellite internet link: the user's dish communicates with a geostationary satellite, which relays traffic to the ISP's ground station connected to the internet.}
\end{popfigure}

The great advantage of satellite is \emph{coverage}: a single satellite can serve an entire country, including the most remote rural zones. The great drawback is \cle{latency}: a geostationary satellite orbits about 36\,000 km above the Earth, so the signal must travel roughly 72\,000 km up and down, adding a noticeable delay (often 500 ms or more). Web browsing still works, but real-time applications such as video calls or online games suffer. Equipment and subscriptions are also relatively expensive.

\courssub{Wireless Local Access: Wi-Fi and Fixed Wireless}

\cle{Wi-Fi} is a \emph{local} wireless technology: a wireless router connected to a wired line (ADSL, fibre) shares the connection over radio waves within a home, school or cybercafé, typically over a few tens of metres. A public Wi-Fi access point is called a \cle{hotspot} (found in airports, hotels, some university campuses).

\cle{Fixed wireless} extends the same idea over longer distances: an ISP installs an antenna on a mast and the customer has a small antenna on the roof, giving broadband without laying cables to each building.

\begin{attention}[Wi-Fi is not the internet connection]
A very common mistake is to confuse \cle{Wi-Fi} with the internet connection itself. Wi-Fi is only the \emph{last few metres}: the wireless link between your device and the router. The router itself must still be connected to the internet by ADSL, fibre, 4G or another technology. A laptop can show ``connected to Wi-Fi'' and still have \emph{no internet access} if the router's own connection is down.
\end{attention}

\courssub{Mobile Networks: from 2G to 5G}

Mobile networks bring internet access through the same cellular masts that carry phone calls. Each generation has brought a leap in speed and services:

\begin{center}
\begin{tabularx}{\linewidth}{|l|l|X|X|}
\hline
\rowcolor{popblueL} \textbf{Generation} & \textbf{Technology} & \textbf{Typical speed} & \textbf{Main services} \\
\hline
2G & GPRS / EDGE & up to about 0.2 Mbps & Text-based browsing, e-mail, basic mobile money \\
\hline
3G & UMTS / HSPA & 1--10 Mbps & Web browsing, social media, small downloads \\
\hline
4G / LTE & LTE & 10--100 Mbps & HD video streaming, video calls, fast downloads \\
\hline
5G & NR & 100 Mbps--1 Gbps+ & Very high speed, very low latency, connected objects \\
\hline
\end{tabularx}
\end{center}

In Cameroon, mobile internet is by far the \emph{dominant} way people get online: MTN and Orange provide 3G and 4G coverage in towns and along main roads, and most Cameroonians experience the internet through a smartphone. Mobile money services (MTN MoMo, Orange Money) run over these same networks. 5G is at an early stage of deployment, concentrated in the largest cities.

\courssub{Comparing the Technologies}

The chart below compares typical \emph{maximum} download speeds on a logarithmic scale, because the values differ enormously.

\begin{popfigure}
\begin{center}
\begin{tikzpicture}
\begin{axis}[
    ybar, bar width=0.9cm,
    width=11.5cm, height=6.5cm,
    ymode=log, log basis y=10,
    ymin=0.01, ymax=3000,
    ylabel={Typical maximum speed (Mbps)},
    symbolic x coords={Dial-up, ADSL, 4G/LTE, Fibre},
    xtick=data,
    nodes near coords,
    every node near coord/.append style={font=\small},
    ymajorgrids, grid style={popline},
    bar shift=0pt,
    fill=popblue, draw=popdark,
]
\addplot coordinates {(Dial-up,0.056) (ADSL,24) (4G/LTE,100) (Fibre,1000)};
\end{axis}
\end{tikzpicture}
\end{center}
\legende{Typical maximum download speeds (logarithmic scale): dial-up 0.056 Mbps, ADSL up to 24 Mbps, 4G up to about 100 Mbps, fibre up to 1000 Mbps (1 Gbps).}
\end{popfigure}

Speed is not the only criterion. A good comparison uses five criteria: \cle{bandwidth}, \cle{availability}, \cle{cost}, \cle{latency} and \cle{mobility}.

\begin{center}
\begin{tabularx}{\linewidth}{|l|X|X|X|}
\hline
\rowcolor{popblueL} \textbf{Technology} & \textbf{Advantage} & \textbf{Disadvantage} & \textbf{Best suited for} \\
\hline
Dial-up & Works on any phone line & Extremely slow; blocks the line & (Historical interest only) \\
\hline
ADSL & Cheap; uses existing lines & Speed falls with distance & Homes near an exchange \\
\hline
Fibre (FTTH) & Highest speed, very reliable & Costly to deploy; limited coverage & Cities, businesses, schools \\
\hline
Satellite & Covers remote areas & High latency; expensive & Rural and isolated zones \\
\hline
Mobile 3G/4G & Mobility; wide coverage & Data bundles costly; coverage gaps & Smartphones, users on the move \\
\hline
\end{tabularx}
\end{center}

\begin{methode}[Choosing an internet connection]
\begin{enumerate}
\item \textbf{Identify the need}: required speed (simple e-mail or video streaming?), number of users, mobility (fixed office or on the move?).
\item \textbf{Check availability} at the location: is there fibre in the quarter? 4G coverage? Only satellite?
\item \textbf{Compare costs}: equipment, installation, and monthly subscription or data bundles.
\item \textbf{Match the technology to the need}: e.g.\ fibre for a city cybercafé, 4G router for a mobile worker, satellite for a rural health centre.
\item \textbf{Verify with a real test} (speed test) before committing to a long contract.
\end{enumerate}
\end{methode}

\begin{exoresolu}[Choosing a connection for a school]
A secondary school in a rural district of the North West Region wants internet for its computer laboratory of 40 machines. There is no telephone line and no fibre; 3G coverage exists but is weak and unstable. Propose a solution and justify it.
\tcblower
\textbf{Solution.} Apply the method: (1) \emph{Need}: 40 simultaneous users require high bandwidth and stability; the school is fixed, so mobility is irrelevant. (2) \emph{Availability}: no ADSL, no fibre; 3G is present but weak. (3) \emph{Cost}: satellite equipment is expensive but is a one-off investment plus subscription. (4) \emph{Match}: a \textbf{satellite link} (dish + router) distributing the connection to the laboratory through Wi-Fi or a wired network is the most suitable option, because it is the only technology offering stable, adequate bandwidth at this location. The weak 3G can serve as a backup. The known drawback of satellite --- latency --- is acceptable for teaching (browsing, e-mail, downloads), though less ideal for video conferencing.
\end{exoresolu}

\begin{exercice}[Comparing two offers]
A trader in Douala hesitates between a 4G router (advertised up to 100 Mbps, paid per data bundle) and an FTTH fibre subscription (up to 50 Mbps, flat monthly rate). She uses the internet heavily all day for her online shop. Advise her, giving one argument based on cost and one based on quality of service.
\end{exercice}

\begin{corrige}[Exercise 1]
\textbf{Cost argument:} with heavy daily use, paying per data bundle on 4G quickly becomes more expensive than a flat-rate fibre subscription, which allows unlimited (or very large) usage for a fixed monthly price. \textbf{Quality argument:} although 4G advertises a higher peak speed (100 Mbps vs 50 Mbps), fibre delivers a more stable, consistent speed with lower latency and no dependence on mobile network congestion --- important for a shop that must stay online all day. Conclusion: the FTTH subscription is the better choice for her fixed, intensive usage; the 4G router could remain as a mobile backup.
\end{corrige}

\begin{attention}[Exam tip]
For the GCE, learn at least \textbf{one advantage and one disadvantage of each} connectivity technology (dial-up, ADSL, cable/fibre, satellite, Wi-Fi, mobile 2G--5G). Questions frequently ask you to ``compare two technologies'' or ``justify a choice for a given situation'' --- always argue from the criteria: \cle{bandwidth}, \cle{availability}, \cle{cost}, \cle{latency} and \cle{mobility}.
\end{attention}

\begin{savaistu}[Internet in Cameroon]
Cameroon illustrates perfectly why several technologies must coexist. Mobile internet (3G/4G) dominates because smartphones are the main device and mobile masts reach far beyond the wired network. Fibre concentrates in Yaoundé and Douala, where CAMTEL, MTN and Orange compete on FTTH offers. And for rural schools, health centres and administrative posts far from any cable, satellite remains the technology of last resort --- which is why national digital-inclusion programmes study all three solutions together.
\end{savaistu}

\begin{aretenir}[Key points]
\begin{itemize}
\item Connections are classified as \cle{dial-up}, \cle{broadband wired}, \cle{wireless} and \cle{mobile}.
\item \cle{Bandwidth} (in bps, kbps, Mbps, Gbps) measures the maximum data rate; \cle{latency} measures the round-trip delay.
\item Dial-up (56 kbps) is historically important but obsolete; ADSL reuses the telephone line with a splitter and offers asymmetric broadband (8--24 Mbps).
\item Fibre optics (FTTH) carries data as light pulses: the highest speeds (up to 1 Gbps) and lowest interference.
\item Satellite covers remote areas but suffers from high latency (about 500 ms).
\item Wi-Fi is only the local wireless link, \emph{not} the internet connection itself.
\item Mobile networks evolved from 2G (GPRS/EDGE) to 5G, multiplying speeds by thousands; in Cameroon, mobile internet dominates.
\item Compare technologies using bandwidth, availability, cost, latency and mobility.
\end{itemize}
\end{aretenir}

\coursec{Technical Requirements, Addressing and Security for Internet Access}

In the previous section we examined the \emph{physical} ways a home, school or business in Cameroon can reach the Internet: ADSL over telephone lines, fibre to the home, 4G/5G mobile networks, satellite links, and the media (copper, fibre, radio) that carry the signals. But owning a cable or a SIM card is not enough. Between your computer and a web server in Douala or abroad, a whole chain of \emph{equipment}, \emph{software}, \emph{addresses} and \emph{protocols} must work together --- and the whole chain must be \emph{protected}. This final section completes the picture: we study what a host needs to connect, how machines are identified on the Internet, which protocols make communication possible, how to manage traffic quality, and how to keep the connection fast and safe. We close with the Cameroonian regulatory framework that governs Internet service provision.

\courssub{Hardware and Software Requirements}

\textbf{Observation.} A student in Buea buys a laptop and a 4G subscription, yet cannot connect her old desktop computer at home. Why? The laptop has a built-in wireless network card; the old desktop has none. Connecting to the Internet begins with the right \cle{hardware}.

\begin{definition}[Hardware requirements for Internet access]
The essential hardware for connecting a host to the Internet comprises:
\begin{itemize}
\item a \cle{Network Interface Card} (NIC): the circuit (Ethernet, Wi-Fi, or mobile broadband) that physically connects the computer to a network medium and handles frames at the data link layer;
\item a \cle{modem} (modulator--demodulator) or \cle{ONT} (Optical Network Terminal): converts digital data into signals suitable for the transmission line (telephone line for ADSL, radio for 4G/5G, light for fibre via an ONT);
\item a \cle{router}: forwards packets between the local network (LAN) and the Internet (WAN), and usually combines modem/ONT interface, switch, Wi-Fi access point and firewall in one box (the ``box'' supplied by CAMTEL, MTN or Orange);
\item \cle{cabling}: UTP (Unshielded Twisted Pair, e.g.\ Cat5e/Cat6) for Ethernet up to \SI{100}{m}, or optical fibre for long distances and very high speeds;
\item a \cle{power backup} (UPS, inverter or generator): essential in Cameroon where power cuts can interrupt connectivity and damage equipment.
\end{itemize}
\end{definition}

\textbf{Software requirements.} Hardware alone is mute; it needs software to speak the language of the Internet.

\begin{itemize}
\item The \cle{operating system network stack}: every modern OS (Windows, Linux, macOS, Android) implements the TCP/IP model --- drivers for the NIC, then IP, then TCP/UDP, then applications.
\item \cle{TCP/IP configuration}: each host needs an IP address, a subnet mask, a default gateway and DNS server addresses. These can be assigned \emph{automatically} by \cle{DHCP} (Dynamic Host Configuration Protocol) or typed in \emph{manually} (static addressing).
\item The \cle{DNS client} (resolver): built into the OS, it translates names such as \texttt{www.minesec.gov.cm} into IP addresses.
\end{itemize}

\begin{attention}[Static or dynamic?]
A common mistake in school computer labs is to assign a static IP address that is already used by another machine or that lies inside the DHCP server's pool. The result is an \emph{IP address conflict}: both machines lose connectivity intermittently. Reserve static addresses for servers and printers, and keep them outside the DHCP range.
\end{attention}

\courssub{IP Addressing: IPv4}

\textbf{Intuition.} To send a letter you need the recipient's address; to send a packet you need the recipient's \cle{IP address}. IPv4, designed in the early 1980s, writes each address as 32 bits, presented as four decimal numbers separated by dots, e.g.\ \texttt{192.168.1.10}. Each number is one octet ($0$--$255$).

\begin{definition}[IPv4 address and subnet mask]
An \cle{IPv4 address} is a 32-bit identifier split by a \cle{subnet mask} into a \emph{network part} and a \emph{host part}. With mask \texttt{255.255.255.0} (or prefix \texttt{/24}), the first 24 bits identify the network and the last 8 bits identify hosts, giving $2^8 - 2 = 254$ usable host addresses (the all-zeros address names the network, the all-ones address is the broadcast).
\end{definition}

\textbf{Private versus public addresses.} Because $2^{32} \approx 4.3$ billion addresses are not enough for the world's devices, three ranges are \emph{private} (usable freely inside organisations, never routed on the public Internet): \texttt{10.0.0.0/8}, \texttt{172.16.0.0/12} and \texttt{192.168.0.0/16}. Your home router gives your phone a private address such as \texttt{192.168.1.10}; on the Internet, all your devices appear under one \emph{public} address owned by your ISP.

\begin{definition}[NAT]
\cle{NAT} (Network Address Translation) is a router mechanism that rewrites the source address and port of outgoing packets so that \emph{many devices with private addresses share a single public IP address}. The router keeps a translation table (private IP + port $\leftrightarrow$ public IP + port) so that replies return to the correct internal device.
\end{definition}

\begin{exemplebox}[NAT at work in a Yaoundé household]
Three phones (\texttt{192.168.1.10}, \texttt{.11}, \texttt{.12}) browse the web through a router whose public address is \texttt{102.244.x.x} (assigned by the ISP). The router replaces each private source address with the public one, using a different source port for each phone. When the web server replies, the router consults its table and delivers the page to the right phone. NAT also hides the internal structure of the network, a modest security side-effect.
\end{exemplebox}

\begin{methode}[Subnetting an IPv4 network (examinable procedure)]
\textbf{Goal:} Divide a given network into $N$ subnets.
\begin{enumerate}
\item Determine the number of bits $b$ to borrow from the host part: smallest $b$ such that $2^b \ge N$.
\item New prefix length $= \text{original prefix} + b$. New mask: set the borrowed bits to 1.
\item Number of hosts per subnet $= 2^{\text{remaining host bits}} - 2$.
\item List subnets: increment the network address by the block size (value of the last borrowed bit).
\item For each subnet: network address (first), broadcast address (last), usable range (network+1 to broadcast-1).
\end{enumerate}
\end{methode}

\begin{exoresolu}[Subnetting a school network]
A school in Bamenda receives the network \texttt{192.168.10.0/24}. It needs subnets for four departments (Administration, Staff room, Computer lab, Students' Wi-Fi). Propose a subnetting plan and give, for the lab subnet, the mask, the network address, the broadcast address and the usable range.
\tcblower
We need at least 4 subnets, so we borrow $b=2$ bits ($2^2 = 4$ subnets). The new prefix is $24 + 2 = 26$, i.e.\ mask \texttt{255.255.255.192} (binary \texttt{11111111.11111111.11111111.11000000}). Each subnet has $2^6 = 64$ addresses, of which $62$ are usable.
\begin{itemize}
\item Administration: \texttt{192.168.10.0/26}, usable \texttt{.1}--\texttt{.62}, broadcast \texttt{.63};
\item Staff room: \texttt{192.168.10.64/26}, usable \texttt{.65}--\texttt{.126}, broadcast \texttt{.127};
\item Computer lab: \texttt{192.168.10.128/26}, usable \texttt{.129}--\texttt{.190}, broadcast \texttt{.191};
\item Students' Wi-Fi: \texttt{192.168.10.192/26}, usable \texttt{.193}--\texttt{.254}, broadcast \texttt{.255}.
\end{itemize}
For the lab: mask \texttt{255.255.255.192}, network \texttt{192.168.10.128}, broadcast \texttt{192.168.10.191}, usable hosts \texttt{192.168.10.129} to \texttt{192.168.10.190} --- enough for 62 computers.
\end{exoresolu}

\courssub{IPv6: The New Generation}

\textbf{Observation.} With billions of phones, vehicles and sensors, IPv4 addresses ran out: the last blocks were distributed to regional registries in the 2010s. The long-term answer is \cle{IPv6}.

\begin{definition}[IPv6 address]
An \cle{IPv6 address} is a 128-bit identifier written as eight groups of four hexadecimal digits separated by colons, e.g.\ \texttt{2001:0db8:0000:0000:0000:ff00:0042:8329}. Leading zeros in a group may be omitted, and \emph{one} run of consecutive all-zero groups may be compressed with \texttt{::}, giving \texttt{2001:db8::ff00:42:8329}.
\end{definition}

\begin{propriete}[Size of the IPv6 address space]
IPv6 provides $2^{128} \approx 3.4 \times 10^{38}$ addresses --- about $7.9 \times 10^{28}$ times more than IPv4. Every device can therefore own a \emph{globally unique} public address, which \emph{eliminates the need for NAT} and restores true end-to-end connectivity. (Statement to know; no proof is examinable.)
\end{propriete}

\textbf{IPv6 address types.}
\begin{itemize}
\item \emph{Unicast} (one interface): includes \emph{global unicast} (routable on Internet, prefix \texttt{2000::/3}), \emph{link-local} (valid only on the local link, prefix \texttt{fe80::/10}, auto-configured), and \emph{unique local} (private, prefix \texttt{fc00::/7}).
\item \emph{Multicast} (one-to-many, prefix \texttt{ff00::/8}) replaces broadcast.
\item \emph{Anycast} (one-to-nearest): same address assigned to multiple interfaces; packet goes to the nearest.
\end{itemize}
Note: IPv6 has \emph{no broadcast} address.

\textbf{Transition mechanisms.} Since the Internet cannot switch overnight, IPv4 and IPv6 coexist through:
\begin{itemize}
\item \cle{Dual stack}: hosts and routers run both protocols simultaneously (most common).
\item \cle{Tunnelling}: IPv6 packets encapsulated inside IPv4 packets (e.g.\ 6to4, Teredo, ISATAP).
\item \cle{Translation}: NAT64/DNS64, converting between the two formats for IPv6-only clients accessing IPv4 servers.
\end{itemize}

\begin{popfigure}
\begin{center}
\begin{tikzpicture}[font=\small, scale=0.95]
% IPv4 box
\draw[fill=popblueL, draw=popblue, thick, rounded corners] (0,0) rectangle (6.4,5.0);
\node[font=\bfseries, text=popdark] at (3.2,4.6) {IPv4};
\node[align=left, text width=5.8cm] at (3.2,2.3) {
\textbf{Size:} 32 bits (4 octets)\\[2pt]
\textbf{Notation:} dotted decimal\\ \hspace{0.4cm}\texttt{192.168.1.10}\\[2pt]
\textbf{Addresses:} $2^{32} \approx 4.3 \times 10^{9}$\\[2pt]
\textbf{Header:} 13 fields, variable\\ \hspace{0.4cm}length (20--60 bytes)\\[2pt]
\textbf{Broadcast:} yes\\[2pt]
\textbf{Config:} DHCPv4 or manual\\[2pt]
\textbf{NAT:} indispensable};
% IPv6 box
\draw[fill=popgreenL, draw=popgreen, thick, rounded corners] (7.0,0) rectangle (13.4,5.0);
\node[font=\bfseries, text=popdark] at (10.2,4.6) {IPv6};
\node[align=left, text width=5.8cm] at (10.2,2.3) {
\textbf{Size:} 128 bits (16 octets)\\[2pt]
\textbf{Notation:} hex + colons\\ \hspace{0.4cm}\texttt{2001:db8::ff00:42:8329}\\[2pt]
\textbf{Addresses:} $2^{128} \approx 3.4 \times 10^{38}$\\[2pt]
\textbf{Header:} 8 fields, fixed\\ \hspace{0.4cm}40 bytes (faster routing)\\[2pt]
\textbf{Broadcast:} no (multicast)\\[2pt]
\textbf{Config:} SLAAC or DHCPv6\\[2pt]
\textbf{NAT:} not needed};
\draw[<->, thick, poporange] (6.4,2.5) -- (7.0,2.5);
\node[text=poporange, font=\footnotesize\bfseries, align=center] at (6.7,2.9) {Coexistence\\(Dual Stack,\\Tunnelling,\\Translation)};
\end{tikzpicture}
\end{center}
\legende{Comparison of IPv4 and IPv6: address space, notation, header and transition.}
\end{popfigure}

\begin{methode}[Configuring a host: static IPv4 and DHCPv6]
\textbf{On Windows (static IPv4):}
\begin{enumerate}
\item Open \emph{Settings $\rightarrow$ Network \& Internet $\rightarrow$ Advanced network settings $\rightarrow$ Change adapter options}, right-click the interface, choose \emph{Properties}.
\item Select \emph{Internet Protocol Version 4 (TCP/IPv4)} $\rightarrow$ \emph{Properties} $\rightarrow$ \emph{Use the following IP address}.
\item Enter the IP address (e.g.\ \texttt{192.168.10.130}), mask (\texttt{255.255.255.192}), default gateway (the router, e.g.\ \texttt{192.168.10.129}) and DNS servers; validate and test with \texttt{ping} and a browser.
\end{enumerate}
\textbf{On Linux (static IPv4, command line, non-persistent):} \texttt{sudo ip addr add 192.168.10.130/26 dev eth0} then \texttt{sudo ip route add default via 192.168.10.129}, and write the DNS server in \texttt{/etc/resolv.conf} (or use \texttt{nmcli} / Netplan for persistence).
\textbf{DHCPv6 client:} on Windows, in \emph{Internet Protocol Version 6 (TCP/IPv6)} properties, choose \emph{Obtain an IPv6 address automatically} and \emph{Obtain DNS server address automatically}. On Linux, enable DHCPv6 in NetworkManager (\texttt{nmcli con mod id "Wired connection 1" ipv6.method dhcp}) or run \texttt{sudo dhclient -6 eth0}. Verify with \texttt{ipconfig /all} (Windows) or \texttt{ip -6 addr} (Linux).
\end{methode}

\courssub{Key Protocols for Internet Access}

\begin{center}
\begin{tabularx}{\linewidth}{|l|X|l|}\hline
\rowcolor{popblueL}\textbf{Protocol} & \textbf{Role} & \textbf{Port / Layer}\\ \hline
HTTP & Transfers web pages in plain text (no confidentiality) & 80 / Application\\ \hline
HTTPS & HTTP encrypted with TLS/SSL: confidentiality, integrity, server authentication (padlock) & 443 / Application\\ \hline
DNS & Translates domain names into IP addresses (and reverse) & 53 / Application\\ \hline
DHCPv4 / DHCPv6 & Automatically assigns IP address, mask, gateway, DNS to hosts & 67/68 (v4), 546/547 (v6) / Application\\ \hline
PPPoE & PPP over Ethernet: authenticates a subscriber (login/password) on an ISP access link, common on ADSL/FTTH & --- / Data Link\\ \hline
VPN & Creates an encrypted tunnel through the public Internet (e.g.\ IPsec, OpenVPN, WireGuard) & Various / Network/Application\\ \hline
TLS/SSL & Cryptographic layer securing HTTP, SMTP, IMAP, etc.\ (SSL is the obsolete ancestor of TLS) & Transport\\ \hline
\end{tabularx}
\end{center}

\begin{savaistu}[The padlock and mobile money]
When you pay a bill by mobile money (MTN MoMo, Orange Money) or check examination results on \texttt{www.gceboard.cm} in Cameroon, HTTPS/TLS encrypts the exchange so that an eavesdropper on the same Wi-Fi network cannot read your PIN or password. ANTIC, the national agency for ICT security, regularly reminds users to check for the padlock and the correct domain name before entering credentials.
\end{savaistu}

\courssub{Quality of Service (QoS)}

\textbf{Observation.} During a video call, if someone in the house starts a large download, the voice becomes choppy. Web pages can wait half a second; a voice packet cannot. \cle{QoS} mechanisms manage this.

\begin{definition}[Quality of Service]
\cle{QoS} is the set of techniques that guarantee measurable performance (bandwidth, delay, jitter, packet loss) to different traffic classes. Two main tools:
\begin{itemize}
\item \cle{Traffic shaping} (or policing): smoothing and limiting the rate of traffic flows so that one user or application cannot saturate the link.
\item \cle{Prioritisation} (queueing): classifying packets into queues (e.g.\ using DSCP -- Differentiated Services Code Point in the IP header) and serving delay-sensitive traffic (VoIP, video conferencing, gaming) \emph{before} elastic traffic (downloads, e-mail, updates).
\end{itemize}
\end{definition}

Routers and switches mark packets and apply queueing disciplines (e.g.\ HTB, FQ\_CoDel). ISPs also use QoS to differentiate subscription offers (e.g.\ guaranteed minimum bandwidth).

\courssub{Security Basics}

\begin{definition}[Essential protection mechanisms]
\begin{itemize}
\item A \cle{firewall} (packet filter / stateful inspection / next-generation) filters incoming and outgoing traffic according to rules (by address, port, protocol, application, user). It may be a dedicated appliance, a router function, or host-based software.
\item An \cle{IDS} (Intrusion Detection System) monitors traffic and logs suspicious patterns; an \cle{IPS} (Intrusion Prevention System) additionally blocks them in real time.
\item \cle{Secure Wi-Fi}: WPA3 (SAE -- Simultaneous Authentication of Equals) encrypts the radio link with strong authentication and forward secrecy; WPA2 (AES-CCMP) remains acceptable. The obsolete WEP and WPA (TKIP) are trivially broken and must never be used.
\item \cle{Phishing awareness}: most intrusions start with a fraudulent e-mail or SMS imitating a bank, a ministry (MINESUP, MINFI) or a mobile-money operator, inviting the victim to a fake page. Always verify the address (HTTPS, correct spelling), never share a PIN or one-time code (OTP).
\item \cle{Updates}: keep OS, browser, router firmware and applications patched against known vulnerabilities.
\end{itemize}
\end{definition}

\begin{attention}[Default passwords and open Wi-Fi]
Routers are delivered with default administrator credentials (often \texttt{admin}/\texttt{admin} or printed on a label) that are publicly known. Keeping them, or running an \emph{open} (unencrypted) Wi-Fi network, lets anyone nearby reconfigure the router, intercept traffic or use the connection for illegal acts --- for which the subscriber may be held legally responsible. First actions on any new router: change the admin password, enable WPA2/WPA3 with a strong passphrase (12+ characters), disable WPS, and update the firmware.
\end{attention}

\courssub{The Cameroonian Regulatory Framework}

In Cameroon, the electronic communications sector is regulated by the \cle{ART} (\emph{Agence de Régulation des Télécommunications}), established by Law No.~98/014 of 14 July 1998 and reorganised by Law No.~2010/013 of 21 December 2010. Its main missions relevant to the Internet are:

\begin{itemize}
\item \cle{Licensing of ISPs}: any operator wishing to provide Internet access (CAMTEL, MTN Cameroon, Orange Cameroon, Nexttel, and others) must obtain a licence or authorisation from ART and comply with its conditions (coverage, quality of service, tariffs, interconnection).
\item \cle{Universal service obligations}: operators contribute to a Universal Service Access Fund intended to extend connectivity to rural and underserved areas, so that Internet access is not limited to Douala and Yaoundé.
\item \cle{Consumer protection and quality control}: ART monitors advertised versus actual quality of service (throughput, latency, availability) and can sanction operators for non-compliance.
\item \cle{Management of national resources}: the \texttt{.cm} country-code top-level domain (ccTLD) and numbering resources fall under national regulation; cybersecurity and critical infrastructure protection are supervised by \cle{ANTIC} (Agence Nationale des Technologies de l'Information et de la Communication), created by Decree No.~2002/053 of 26 March 2002 and reinforced by Law No.~2010/012 on cybersecurity and cybercrime.
\end{itemize}

\begin{aretenir}[Key points of Section 4]
\begin{itemize}
\item Connecting a host requires hardware (NIC, modem/ONT, router, cabling, power backup) and software (TCP/IP stack, DHCP/DNS client).
\item IPv4: 32 bits, dotted decimal; private ranges + NAT allow many devices to share one public address. Subnetting borrows host bits to create subnets (examinable calculation).
\item IPv6: 128 bits, hexadecimal, $2^{128}$ addresses, no NAT, no broadcast (multicast instead); coexistence via dual stack, tunnelling, translation.
\item Essential protocols: HTTP/HTTPS, DNS, DHCP, PPPoE, VPN, TLS/SSL.
\item QoS (shaping, prioritisation/DSCP) protects real-time traffic; security rests on firewalls, IDS/IPS, WPA3, updates and user vigilance against phishing.
\item In Cameroon, ART licenses ISPs and imposes universal service obligations; ANTIC oversees ICT security and the \texttt{.cm} domain.
\end{itemize}
\end{aretenir}

\begin{exercice}[Addresses, protocols and security]
\begin{enumerate}
\item A host has address \texttt{172.20.5.44} with mask \texttt{255.255.0.0}. Is this address private or public? Justify using the RFC 1918 ranges.
\item Compress the IPv6 address \texttt{2001:0db8:0000:0000:0000:0000:0000:0001} following the standard rules.
\item Explain why NAT complicates hosting a public web server at home, and how IPv6 solves the problem without compromising security.
\item List four concrete measures to secure the Wi-Fi network of a cybercafé in Douala.
\item A small business has the network \texttt{192.168.50.0/24}. It needs 6 subnets. Determine the new mask, the number of usable hosts per subnet, and the network/broadcast addresses of the 3rd subnet.
\end{enumerate}
\end{exercice}

\begin{corrige}[Exercise 1]
\begin{enumerate}
\item \texttt{172.20.5.44} belongs to the range \texttt{172.16.0.0}--\texttt{172.31.255.255} (the \texttt{172.16.0.0/12} private block defined in RFC 1918), so it is a \emph{private} address, not routable on the public Internet.
\item \texttt{2001:db8::1} (leading zeros removed in each group; the five consecutive all-zero groups compressed with a single \texttt{::}).
\item With NAT, internal machines have private addresses unreachable from outside; hosting a server requires manual \emph{port forwarding} (DNAT) on the router, which is fragile, limited to one server per port, and breaks end-to-end connectivity. IPv6 gives every device a global unicast address, so the server is directly reachable from the Internet; access is controlled by a stateful firewall on the router (allow inbound TCP/80,443 to that specific IPv6 address) --- security is maintained without address translation.
\item Enable WPA2/WPA3 (AES) with a strong passphrase (changed periodically); change the router's default administrator password; update firmware regularly; disable WPS; enable client isolation (AP isolation) so clients cannot talk to each other; optionally use a captive portal with RADIUS authentication for accountability.
\item Need 6 subnets $\rightarrow$ borrow $b=3$ bits ($2^3=8$ subnets). New prefix $= 24+3 = 27$. Mask \texttt{255.255.255.224}. Block size $= 32$ addresses. Usable hosts $= 32-2 = 30$.
Subnets: 0: \texttt{.0}, 1: \texttt{.32}, 2: \texttt{.64}, \textbf{3: \texttt{.96}}, 4: \texttt{.128}, 5: \texttt{.160}, 6: \texttt{.192}, 7: \texttt{.224}.
3rd subnet (index 2 if 0-based, but usually "3rd" means index 2 or 3? Standard wording: 1st=.0, 2nd=.32, 3rd=.64). Let's assume 1st = .0/27. 3rd = .64/27.
Network: \texttt{192.168.50.64}, Broadcast: \texttt{192.168.50.95}, Usable: \texttt{.65}--\texttt{.94}.
\end{enumerate}
\end{corrige}

\newpage

\coursec{Integration activity}

\courssub{Aims of the activity}

This integration activity closes the chapter \emph{Reviewing the History of the Internet}. It asks you to act as an ICT consultant: you will use everything learned about the \cle{history of the Internet} (Lesson 36) and about \cle{connectivity technologies and requirements} (Lesson 37) to solve a realistic Cameroonian problem. By the end of the activity, you should be able to:

\begin{itemize}
\item mobilise your knowledge of the \cle{historical development} of the Internet (from ARPANET to the modern commercial Internet) to understand why certain technologies dominate today;
\item compare \cle{Internet connectivity technologies} (fibre optic, 4G mobile broadband, satellite, ADSL) and justify a choice in terms of bandwidth, latency, availability and cost;
\item calculate the \cle{bandwidth requirements} of a real institution from usage scenarios;
\item design a simple \cle{Local Area Network} (LAN) with an Internet gateway, including redundancy (backup link);
\item propose an \cle{IPv4/IPv6 addressing plan} and \cle{basic security measures} (firewall, authentication, filtering);
\item communicate a technical solution clearly in a short oral presentation, as an ICT professional would.
\end{itemize}

The activity is done in \textbf{groups of 4} and is assessed with a rubric covering \emph{technical justification} (40\%), \emph{bandwidth and cost analysis} (30\%), and \emph{clarity of the diagram and the 5-minute pitch} (30\%).

\courssub{Situation}

\textbf{Connecting Collège Bilingue de Mandjou (East Region)}\par

Collège Bilingue de Mandjou is a government secondary school located on the outskirts of Bertoua, capital of the East Region of Cameroon. The school has \textbf{300 students} and \textbf{20 teachers}. Thanks to a donation, it now owns a \textbf{computer laboratory of 30 PCs}, but none of them is connected to the Internet. The principal wants to use the lab for ICT practical lessons, online research, and \emph{simultaneous video lessons} streamed from the Ministry's distance-education platform, especially for subjects where qualified teachers are scarce.

A feasibility visit reveals the following facts:

\begin{itemize}
\item A \textbf{CAMTEL fibre optic backbone} passes along the main road, about \SI{800}{m} from the school. CAMTEL quotes a one-off installation fee of CFA 1\,500\,000 and a monthly subscription of CFA 150\,000 for a dedicated \SI{100}{Mbps} link.
\item \textbf{MTN 4G} coverage in the area is good (measured average: \SI{20}{Mbps} download, \SI{5}{Mbps} upload, latency about \SI{60}{ms}). A 4G router costs CFA 60\,000; data bundles cost about CFA 10\,000 for 50\,GB.
\item \textbf{Satellite Internet} (VSAT) is available everywhere but costs CFA 2\,500\,000 for equipment plus CFA 200\,000 per month, with a high latency of about \SI{600}{ms} and a fair-use data cap.
\item The school's total budget for the project (equipment + installation + one year of subscription) is \textbf{CFA 5\,000\,000}.
\item The lab is a single hall; the administration block (5 offices) is \SI{40}{m} away. Teachers should also get Wi-Fi access in the staff room.
\item The principal insists: \emph{``If the main link fails during an examination session, we must still be able to go online.''}
\end{itemize}

\begin{attention}[Budget discipline]
The CFA 5\,000\,000 must cover \emph{everything}: the primary link installation, one year of subscription, the backup link, the router, switch, Wi-Fi access points, cabling and a small firewall. A solution that exceeds the budget is rejected, however brilliant it may be technically. Always draw up the full budget table \emph{before} finalising your technical choices.
\end{attention}

\courssub{Tasks}

\textbf{Task 1 — Choosing the primary connectivity technology.}\par
For each of the three options (CAMTEL fibre, MTN 4G, satellite VSAT), identify its strengths and weaknesses in terms of bandwidth, latency, reliability and cost. Then choose the \emph{primary} link and the \emph{backup} link, justifying your choice with at least three technical arguments. Recall from the history of the Internet why fibre has become the dominant fixed-access medium.

\textbf{Task 2 — Calculating the required bandwidth.}\par
The school wants up to \textbf{30 simultaneous video streams} (one per PC), each requiring \SI{2}{Mbps}. Additionally, allow \SI{10}{Mbps} for administration, messaging and background traffic.
\begin{enumerate}
\item Compute the total bandwidth needed.
\item Check whether your chosen primary link satisfies this need, with a safety margin.
\item Estimate the monthly data volume if the lab streams video 4 hours per day, 20 days per month (1 hour of a \SI{2}{Mbps} stream $\approx$ \SI{0.9}{GB}). Comment on whether a 4G-only solution would be viable.
\end{enumerate}

\textbf{Task 3 — Drafting the network diagram.}\par
Draw a labelled diagram of the school network showing: the primary link, the backup link, a router/gateway, a firewall, a central switch in the lab, the 30 PCs, at least two Wi-Fi access points (lab and staff room) and the administration block. Indicate the transmission medium used between buildings and justify it.

\textbf{Task 4 — Addressing and security plan.}\par
\begin{enumerate}
\item Propose a private IPv4 addressing plan (network address, subnet mask, ranges for PCs, servers, Wi-Fi users) using CIDR notation, and explain the role of NAT and DHCP.
\item Explain briefly why the school should also be ready for IPv6.
\item List at least four basic security measures (e.g.\ firewall rules, WPA2/WPA3 Wi-Fi authentication, content filtering, user accounts, updates) and state what threat each one addresses.
\end{enumerate}

\textbf{Task 5 — The 5-minute pitch.}\par
Prepare a 5-minute group presentation for the principal and the Parent-Teacher Association: one slide on the chosen technology, one on the bandwidth calculation, one on the diagram, and one on the budget table. Every group member must speak.

\courssub{Model answer}

\textbf{Task 1 — Choice of technology.}\par
The comparison table summarises the three options. First-year costs are computed as: installation (or equipment) $+\ 12 \times$ monthly subscription.

\begin{center}
\begin{tabularx}{\linewidth}{|l|X|X|X|}
\hline
\rowcolor{popblueL} \textbf{Criterion} & \textbf{CAMTEL fibre} & \textbf{MTN 4G} & \textbf{Satellite VSAT} \\
\hline
Bandwidth & \SI{100}{Mbps} dedicated & $\approx$ \SI{20}{Mbps}, shared & 10--30 Mbps, capped \\
\hline
Latency & Very low ($<$ \SI{10}{ms}) & $\approx$ \SI{60}{ms} & $\approx$ \SI{600}{ms} \\
\hline
Reliability & High, stable & Variable (congestion, signal quality) & Affected by heavy rain \\
\hline
Data limit & Unlimited & Pay per bundle (50\,GB) & Fair-use cap \\
\hline
Cost (year 1) & 1.5\,M $+\ 12\times 150$\,k $=$ 3.3\,M FCFA & Router 60\,k $+$ bundles & 2.5\,M $+\ 12\times 200$\,k $=$ 4.9\,M FCFA \\
\hline
\end{tabularx}
\end{center}

\textbf{Decision:} the \cle{primary link} is the \textbf{CAMTEL fibre}, for three main reasons: (1) it offers \SI{100}{Mbps} \emph{dedicated}, enough for the 30 simultaneous video streams; (2) its latency is very low, which is essential for interactive video lessons; (3) its data is unlimited and its total first-year cost (3.3\,M FCFA) fits the budget. The \cle{backup link} is an \textbf{MTN 4G router}: cheap, quick to activate, and sufficient for emergency text-based work (e-mail, result submission) if the fibre is cut. Satellite is rejected: it is the most expensive option (4.9\,M FCFA in year one, almost the whole budget), has the highest latency (about \SI{600}{ms}, poor for interactive video) and a data cap.

Historically, fibre dominates fixed access because, from the 1990s onwards, the explosive growth of Internet traffic — the Web, then audio and video streaming — demanded capacities that copper telephone lines (ADSL) could not offer, while the cost of optical fibre and optical equipment fell steadily. This is the same evolution that carried the Internet from the low-speed ARPANET links of the 1970s to today's broadband networks.

\textbf{Task 2 — Bandwidth calculation.}\par
The total bandwidth is the sum of the simultaneous needs:
\begin{align*}
\text{Video streams} &= 30 \times \SI{2}{Mbps} = \SI{60}{Mbps}\\
\text{Other traffic} &= \SI{10}{Mbps}\\
\text{Total required} &= 60 + 10 = \SI{70}{Mbps}
\end{align*}
The fibre link offers \SI{100}{Mbps}. The safety margin is
\[
\frac{100 - 70}{100} = 0.30 = 30\%
\]
of spare capacity: the link is \textbf{acceptable}, even if usage grows slightly.

For the data volume, note first where the given figure comes from: a \SI{2}{Mbps} stream transfers $2\ \mathrm{Mbit/s} \div 8 = 0.25\ \mathrm{MB/s}$, i.e.\ $0.25 \times 3600 = 900\ \mathrm{MB} \approx \SI{0.9}{GB}$ per hour. Hence:
\[
30 \text{ PCs} \times 0.9\ \mathrm{GB/h} \times 4\ \mathrm{h/day} \times 20\ \mathrm{days} = \SI{2160}{GB} \approx \SI{2.2}{TB} \text{ per month.}
\]
On 4G at CFA 10\,000 per 50\,GB bundle, this would cost
\[
\frac{2160}{50} \times 10\,000 \approx \text{CFA } 432\,000 \text{ per month,}
\]
nearly three times the fibre subscription, for a slower, shared link with higher latency. A 4G-only solution is therefore \emph{not viable} as the primary link — but it remains excellent as an occasional backup.

\textbf{Task 3 — Network diagram.}

\begin{popfigure}
\begin{tikzpicture}[font=\small, node distance=1.1cm,
box/.style={draw=popblue, thick, rounded corners, fill=popblueL, align=center, text width=2.4cm, minimum height=0.9cm}]
\node[box] (fibre) {CAMTEL Fibre (primary)};
\node[box, right=1.6cm of fibre] (g4) {MTN 4G router (backup)};
\node[box, below=1.2cm of fibre, xshift=1.9cm] (fw) {Router + Firewall (gateway)};
\node[box, below=1.2cm of fw] (sw) {Central switch (lab)};
\node[box, below left=1.2cm and -0.6cm of sw] (pcs) {30 lab PCs (cable)};
\node[box, below=1.2cm of sw] (ap1) {Wi-Fi AP (lab)};
\node[box, below right=1.2cm and -0.6cm of sw] (ap2) {Wi-Fi AP (staff room)};
\node[box, right=1.6cm of sw] (admin) {Admin block (5 offices)};
\draw[thick, popdark] (fibre) -- (fw);
\draw[thick, popdark] (g4) -- (fw);
\draw[thick, popdark] (fw) -- node[right, popdark]{\scriptsize Ethernet} (sw);
\draw[thick, popdark] (sw) -- (pcs);
\draw[thick, popdark] (sw) -- (ap1);
\draw[thick, popdark] (sw) -- (ap2);
\draw[thick, poporange] (sw) -- node[above, poporange]{\scriptsize fibre \SI{40}{m}} (admin);
\end{tikzpicture}
\legende{Proposed network for Collège Bilingue de Mandjou. Both Internet links enter the router/firewall; the switch distributes the connection inside the lab. The inter-building link uses outdoor fibre to avoid interference and signal loss over \SI{40}{m}.}
\end{popfigure}

The two links (fibre and 4G) both connect to the \cle{router/gateway}, which can switch to the backup automatically if the fibre fails. Inside the lab, PCs are wired to the switch with Ethernet cable (cheaper and more stable than Wi-Fi for fixed machines). The link to the administration block, \SI{40}{m} away, uses \textbf{outdoor optical fibre}: copper Ethernet is limited to about \SI{100}{m} and is sensitive to electromagnetic interference and lightning, whereas fibre is immune and future-proof.

\textbf{Task 4 — Addressing and security.}\par
\textbf{IPv4 plan.} We use the private range 192.168.10.0/24 (mask 255.255.255.0, giving $2^8 - 2 = 254$ usable host addresses). The router performs \cle{NAT} and a \cle{DHCP} server assigns addresses automatically.

\begin{center}
\begin{tabularx}{\linewidth}{|l|X|X|}
\hline
\rowcolor{popblueL} \textbf{Element} & \textbf{Value} & \textbf{Remark} \\
\hline
Network address & 192.168.10.0/24 & Mask 255.255.255.0, 254 usable hosts \\
\hline
Router / gateway & 192.168.10.1 & Static \\
\hline
Switch, APs, server & 192.168.10.2 -- 192.168.10.10 & Static (reserved, outside DHCP pool) \\
\hline
Lab PCs (DHCP pool) & 192.168.10.50 -- 192.168.10.100 & 30 PCs + spares \\
\hline
Wi-Fi users (DHCP pool) & 192.168.10.150 -- 192.168.10.220 & Staff and students \\
\hline
\end{tabularx}
\end{center}

\cle{NAT} (Network Address Translation) hides all the private addresses behind the single public address provided by CAMTEL, saving public addresses and adding a first layer of protection. \cle{DHCP} (Dynamic Host Configuration Protocol) automates the configuration of each machine (address, mask, gateway, DNS) and avoids address conflicts and typing errors.

\textbf{IPv6 readiness:} the global pool of IPv4 addresses is exhausted, and ISPs (including Cameroonian operators) progressively deploy IPv6. The school should therefore buy \emph{dual-stack} equipment (supporting IPv4 and IPv6) so that no replacement is needed when IPv6 becomes standard.

\textbf{Security measures} (at least four; each is matched to the threat it addresses):
\begin{enumerate}
\item \textbf{Firewall} on the router: blocks unsolicited incoming connections (threat: intrusions, port scans from the Internet).
\item \textbf{WPA2/WPA3} with a strong password on the Wi-Fi access points (threat: unauthorised access, bandwidth theft, eavesdropping).
\item \textbf{Content filtering} at the gateway (threat: inappropriate sites for students, phishing and malware sites).
\item \textbf{Individual user accounts} with passwords and limited rights on the PCs (threat: data deletion, installation of unwanted software, misuse).
\item \textbf{Regular updates and antivirus} on all PCs (threat: viruses, ransomware exploiting known flaws).
\end{enumerate}

\textbf{Task 5 — Budget summary for the pitch.}

\begin{center}
\begin{tabularx}{\linewidth}{|X|l|}
\hline
\rowcolor{popblueL} \textbf{Item} & \textbf{Cost (FCFA)} \\
\hline
Fibre installation + 1 year subscription ($1.5\,\text{M} + 12 \times 150\,000$) & 3\,300\,000 \\
\hline
4G backup router + 1 year of bundles ($60\,000 + 12 \times 10\,000$) & 180\,000 \\
\hline
Router/firewall, switch, 2 access points, cabling & 1\,200\,000 \\
\hline
Contingency (about 5\%) & 250\,000 \\
\hline
\textbf{Total} & \textbf{4\,930\,000} \\
\hline
\end{tabularx}
\end{center}

The project fits the CFA 5\,000\,000 budget (with CFA 70\,000 to spare), delivers \SI{100}{Mbps} with a 30\% safety margin, includes the backup link demanded by the principal, and secures the network. The pitch concludes: \emph{``Fibre for performance, 4G for safety, security for trust.''}

\begin{attention}[Common mistakes to avoid]
\begin{itemize}
\item Forgetting the backup link, or the one-year subscription, in the budget.
\item Confusing \si{Mbps} (mega\emph{bits} per second, used for links) with \si{MB} (mega\emph{bytes}, used for storage): $1\ \mathrm{byte} = 8\ \mathrm{bits}$, so a \SI{2}{Mbps} stream downloads only \SI{0.25}{MB} per second.
\item Assigning public addresses inside the LAN, omitting the subnet mask, or placing static addresses inside the DHCP pool (risk of conflicts).
\item Drawing a diagram without labelling the transmission media, or without the firewall between the Internet links and the LAN.
\item Choosing a technology on bandwidth alone, without checking latency, data caps and total cost over one year.
\end{itemize}
\end{attention}

\newpage

\coursec{Methods and worked examples}

\coursec{Origins and Early Development of the Internet}

\courssub{Skill 1 — Placing Key Events in the History of the Internet (Timeline Questions)}

\begin{definition}[Internet]
The \cle{Internet} is a global network of interconnected computer networks that communicate using a common set of protocols, the TCP/IP protocol suite, allowing users worldwide to share information and services.
\end{definition}

\begin{definition}[World Wide Web]
The \cle{World Wide Web} (WWW or Web) is an information system of interlinked hypertext documents accessed via the Internet. It uses URLs to identify resources, HTTP to transfer them, and HTML to format them.
\end{definition}

\begin{aretenir}[Internet vs. Web — The Distinction That Earns Marks]
\begin{itemize}
\item The \textbf{Internet} is the \emph{infrastructure} (hardware + protocols): a network of networks born in 1969 (ARPANET) / 1983 (TCP/IP).
\item The \textbf{Web} is a \emph{service} that runs \emph{on top of} the Internet, invented in 1989--1991 by Tim Berners-Lee at CERN.
\item \textbf{Never} write ``the Internet was invented in 1989'' — that date belongs to the Web.
\end{itemize}
\end{aretenir}

\begin{methode}[How to answer a GCE question on the history of the Internet]
\begin{enumerate}
\item \textbf{Identify the period asked.} The GCE typically asks about one of four eras: the pre-Internet era (before 1969), the ARPANET era (1969--1982), the TCP/IP era (1983--1990), and the Web era (1990 to today).
\item \textbf{Recall the anchor dates.} Memorise the short list: 1969 (ARPANET first nodes), 1973--1974 (TCP designed by Cerf and Kahn), 1983 (ARPANET switches to TCP/IP, 1 January), 1989--1991 (Tim Berners-Lee invents the World Wide Web at CERN), 1993 (Mosaic, the first popular graphical browser).
\item \textbf{Distinguish Internet from Web.} The \cle{Internet} is the global network of networks (hardware + protocols); the \cle{World Wide Web} is a service (hypertext pages linked by URLs, transferred with HTTP) that \emph{runs on} the Internet.
\item \textbf{Structure the answer.} For an essay question, use chronological order, give the date, the event, the key people or organisations, and the consequence of each step.
\item \textbf{Link to Cameroon when asked.} Mention that Cameroon connected progressively in the late 1990s, with CAMTEL as the historic operator and ANTIC created later to regulate and secure the digital space.
\end{enumerate}
\end{methode}

\begin{savaistu}[Cameroon's Early Internet]
Cameroon's first Internet connection was established in the late 1990s via CAMTEL. The \cle{ANTIC} (Agence Nationale des Technologies de l'Information et de la Communication) was created in 2002 to regulate, promote and secure ICT usage. Today, the country is served by submarine cables (WACS, SAT-3, ACE) landing at Kribi and Limbe, feeding the national fibre backbone operated by CAMTEL.
\end{savaistu}

\begin{exoresolu}[Timeline of the Internet]
\textbf{Statement (GCE A-Level style).}\\
(a) Define the term \emph{Internet}. \hfill (2 marks)\\
(b) Arrange the following events in chronological order and give the approximate year of each: invention of the World Wide Web; adoption of TCP/IP by ARPANET; creation of ARPANET; release of the Mosaic browser. \hfill (4 marks)\\
(c) State two differences between the Internet and the World Wide Web. \hfill (4 marks)
\tcblower
\textbf{Detailed solution.}\\
(a) \textbf{Definition.} The Internet is a global network of interconnected computer networks that communicate using a common set of protocols, the TCP/IP protocol suite, allowing users worldwide to share information and services. \emph{(1 mark for ``network of networks'', 1 mark for the common protocols / global communication.)}

(b) \textbf{Chronological order.}
\begin{enumerate}
\item \textbf{1969 — Creation of ARPANET.} Funded by the US Department of Defense (ARPA), the first nodes connected four American universities (UCLA, Stanford, UC Santa Barbara, Utah); it used packet switching.
\item \textbf{1983 — Adoption of TCP/IP by ARPANET.} On 1 January 1983 ARPANET switched from the NCP protocol to TCP/IP, designed by Vinton Cerf and Robert Kahn; this date is often considered the ``birth'' of the modern Internet.
\item \textbf{1989--1991 — Invention of the World Wide Web.} Tim Berners-Lee, at CERN in Geneva, proposed and built the Web: URLs, HTTP and HTML, plus the first browser/editor.
\item \textbf{1993 — Release of Mosaic.} The first widely used graphical web browser, developed at NCSA, which made the Web accessible to the general public and triggered its explosive growth.
\end{enumerate}

(c) \textbf{Two differences.}
\begin{center}
\begin{tabularx}{\linewidth}{|l|X|X|}
\hline
\rowcolor{popblueL} \textbf{Criterion} & \textbf{Internet} & \textbf{World Wide Web} \\
\hline
Nature & Physical/logical infrastructure: a network of networks (cables, routers, protocols) & A service: a collection of interlinked hypertext documents accessed via browsers \\
\hline
Birth & 1969 (ARPANET) / 1983 (TCP/IP) & 1989--1991 (Tim Berners-Lee at CERN) \\
\hline
Key technologies & TCP/IP, routers, DNS & HTTP, HTML, URLs, browsers \\
\hline
\end{tabularx}
\end{center}
\emph{Any two well-expressed differences earn the marks. The classic trap — saying the Web and the Internet are the same thing — must be avoided.}
\end{exoresolu}

\courssub{Skill 2 — Explaining Packet Switching and the Role of TCP/IP}

\begin{definition}[Packet Switching]
\cle{Packet switching} is a data transmission technique in which a message is divided into small units called \cle{packets}; each packet carries a header with source and destination addresses; packets travel independently through the network (possibly by different routes) and are reassembled in the correct order at the destination.
\end{definition}

\begin{propriete}[Advantages of Packet Switching over Circuit Switching]
\begin{itemize}
\item \textbf{Efficiency:} transmission links are shared between many users (statistical multiplexing); bandwidth is not wasted during silent periods.
\item \textbf{Robustness:} if a router or link fails, packets are rerouted along another path; the communication survives partial network failures (a design goal of ARPANET).
\item \textbf{No dedicated circuit:} no need to establish and reserve an end-to-end line before communicating, unlike the classical telephone network (PSTN).
\end{itemize}
\end{propriete}

\begin{definition}[TCP and IP]
\begin{itemize}
\item \cle{IP} (Internet Protocol): handles \emph{addressing and routing} of packets. Every host has an IP address; routers forward packets based on the destination IP.
\item \cle{TCP} (Transmission Control Protocol): handles \emph{reliable delivery}. It numbers packets, detects losses, requests retransmission, and reassembles them in order.
\end{itemize}
\end{definition}

\begin{methode}[How to explain why the Internet works: packet switching + TCP/IP]
\begin{enumerate}
\item \textbf{Define packet splitting.} A message is split into small blocks called packets; each packet carries a header with source and destination addresses.
\item \textbf{Explain independent routing.} Packets may travel by different routes through the network; routers read the destination address and forward each packet along the best available path.
\item \textbf{Describe reassembly.} At the destination, TCP checks that all packets arrived, requests retransmission of any lost or corrupted packet, and reorders them.
\item \textbf{Give the advantages.} Efficient use of links (sharing), robustness (rerouting), no dedicated line needed.
\item \textbf{Use an analogy.} Sending a book as several numbered parcels through a postal network: each parcel may take a different road, but the numbering lets the receiver rebuild the book.
\item \textbf{Conclude.} TCP/IP became the standard in 1983 and is the reason networks of different types can intercommunicate — hence ``\emph{inter}-net''.
\end{enumerate}
\end{methode}

\begin{exoresolu}[Packet switching and TCP/IP]
\textbf{Statement.}\\
(a) What is packet switching? \hfill (2 marks)\\
(b) Give two advantages of packet switching over circuit switching. \hfill (2 marks)\\
(c) A school in Bamenda sends a 4 MB photo by email to a school in Buea. Explain, in four steps, what happens to the photo from the moment it leaves the sender's computer until it is displayed on the receiver's screen. \hfill (4 marks)
\tcblower
\textbf{Detailed solution.}\\
(a) \textbf{Packet switching} is a data transmission technique in which a message is divided into small units called packets; each packet is labelled with the destination address, travels independently through the network (possibly by different routes), and all packets are reassembled in the correct order at the destination.

(b) \textbf{Advantages (any two).}
\begin{itemize}
\item \textbf{Efficiency:} transmission links are shared between many users instead of being reserved for a single conversation, so bandwidth is not wasted.
\item \textbf{Robustness:} if a router or link fails, packets are simply rerouted along another path; the communication survives partial network failures.
\item \textbf{No dedicated circuit:} there is no need to establish and reserve an end-to-end line before communicating, unlike the classical telephone network.
\end{itemize}

(c) \textbf{The four steps.}
\begin{enumerate}
\item \textbf{Segmentation:} the sender's computer (TCP layer) splits the 4 MB photo into numbered packets and adds headers containing the destination IP address of the Buea computer.
\item \textbf{Routing:} each packet travels across the network (through the ISP's routers, e.g.\ via CAMTEL's backbone); routers read the destination address and forward each packet along the best available path — packets may take different routes.
\item \textbf{Reassembly and checking:} the receiver's TCP layer checks that all packets arrived, requests retransmission of any lost or corrupted packet, and reorders them.
\item \textbf{Delivery:} the complete photo is reconstructed and handed to the email application, which displays it on the screen.
\end{enumerate}
\end{exoresolu}

\coursec{Expansion, Commercialisation and the Modern Internet}

\courssub{Skill 3 — Identifying and Comparing Internet Connectivity Technologies}

\begin{aretenir}[Main Families of Internet Access]
\begin{itemize}
\item \textbf{Wired:} Dial-up (PSTN), ADSL, Fibre optic (FTTH), Cable (DOCSIS).
\item \textbf{Wireless:} Mobile networks (2G/3G/4G/5G), Wi-Fi (local), Satellite (VSAT, LEO constellations), WiMAX.
\end{itemize}
\end{aretenir}

\begin{methode}[How to choose or compare an Internet access technology]
\begin{enumerate}
\item \textbf{List the main families.} Wired: dial-up (PSTN), ADSL, fibre optic (FTTH), cable. Wireless: mobile networks (2G/3G/4G/5G), Wi-Fi, satellite, WiMAX.
\item \textbf{Compare with fixed criteria.} Always compare using: maximum bandwidth (speed), coverage/availability, cost, mobility, and reliability (sensitivity to weather, distance, congestion).
\item \textbf{Know the orders of magnitude.} Dial-up: up to 56 kbit/s. ADSL: a few Mbit/s (degrades with distance from the exchange). Fibre: tens to hundreds of Mbit/s and more (symmetric possible). 4G: tens of Mbit/s in practice. Satellite: wide coverage but high latency (GEO) or lower latency (LEO).
\item \textbf{Match technology to context.} A rural village without a telephone line $\rightarrow$ mobile network or satellite; a bank in Douala needing high speed and stability $\rightarrow$ fibre; a home user $\rightarrow$ ADSL or 4G router.
\item \textbf{In the exam, justify.} A mark is earned only if the choice is justified with a criterion (speed, cost, coverage), not merely named.
\end{enumerate}
\end{methode}

\begin{attention}[Common Trap: Confusing Local and Access Technologies]
Wi-Fi is a \emph{local} wireless technology (last few metres). It is \textbf{not} an Internet access technology by itself — it connects your device to a router that uses ADSL, fibre, or 4G for the actual Internet link. Do not list ``Wi-Fi'' as an ISP access type in a comparison table.
\end{attention}

\begin{exoresolu}[Choosing a connectivity technology]
\textbf{Statement.}\\
A cybercafé owner in Ngaoundéré asks for advice. The town has a CAMTEL telephone line and good 4G coverage from MTN and Orange, but no fibre yet.\\
(a) Copy and complete the comparison table below. \hfill (4 marks)
\begin{center}
\begin{tabularx}{\linewidth}{|l|X|X|}
\hline
\rowcolor{popblueL} \textbf{Criterion} & \textbf{ADSL} & \textbf{4G (mobile)} \\
\hline
Transmission medium & & \\
\hline
Typical speed & & \\
\hline
One advantage & & \\
\hline
One limitation & & \\
\hline
\end{tabularx}
\end{center}
(b) Recommend one technology for the cybercafé and justify your choice with two arguments. \hfill (3 marks)
\tcblower
\textbf{Detailed solution.}\\
(a) \textbf{Completed table.}
\begin{center}
\begin{tabularx}{\linewidth}{|l|X|X|}
\hline
\rowcolor{popblueL} \textbf{Criterion} & \textbf{ADSL} & \textbf{4G (mobile)} \\
\hline
Transmission medium & Existing copper telephone line (PSTN) & Radio waves via a mobile base station (antenna) \\
\hline
Typical speed & A few Mbit/s (e.g.\ 1--20 Mbit/s), decreasing with distance from the exchange & Tens of Mbit/s in good coverage, shared between users of the cell \\
\hline
One advantage & Stable, dedicated bandwidth; flat-rate (unlimited) subscriptions common & Quick to deploy, no line works needed; usable on the move \\
\hline
One limitation & Requires a telephone line; speed drops far from the exchange & Speed varies with congestion and signal strength; data bundles can be expensive \\
\hline
\end{tabularx}
\end{center}

(b) \textbf{Recommendation.} For a cybercafé, \textbf{ADSL} is recommended (with a 4G router as backup).
\begin{itemize}
\item \textbf{Argument 1 — Cost model:} a cybercafé consumes large volumes of data all day; ADSL flat-rate subscriptions allow unlimited usage, whereas 4G pay-per-gigabyte bundles would be very costly for many simultaneous customers.
\item \textbf{Argument 2 — Stability:} ADSL offers a dedicated line whose bandwidth is not shared with neighbouring users, unlike a 4G cell whose speed collapses at peak hours when many subscribers are connected.
\end{itemize}
\emph{(A reasoned answer choosing 4G — e.g.\ if the telephone line is unreliable — can also earn marks provided two valid criteria are given.)}
\end{exoresolu}

\coursec{Internet Connectivity Technologies and Media}

\courssub{Skill 4 — Listing the Requirements for an Internet Connection}

\begin{aretenir}[Four Categories of Requirements]
To connect a computer to the Internet, you need:
\begin{enumerate}
\item \textbf{Hardware:} computer/smartphone, network interface (Ethernet/Wi-Fi), modem/router (ADSL modem, 4G router, ONT for fibre), transmission medium (telephone line, fibre, SIM card).
\item \textbf{An ISP subscription:} contract with an \cle{Internet Service Provider} (CAMTEL, MTN Cameroon, Orange Cameroon, etc.) supplying the link and an IP address.
\item \textbf{Software:} OS with TCP/IP stack configured, applications (web browser, email client).
\item \textbf{Authentication:} username/password or activated SIM with data plan.
\end{enumerate}
\end{aretenir}

\begin{methode}[How to answer ``State the requirements for connecting a computer to the Internet'']
\begin{enumerate}
\item \textbf{Think in four categories.} (i) Hardware, (ii) an access provider, (iii) software, (iv) a subscription/account.
\item \textbf{Hardware.} A computer (or smartphone), a \cle{modem} or router (ADSL modem, 4G router, ONT for fibre) or a network interface (Ethernet/Wi-Fi card), and the transmission medium (telephone line, fibre, SIM card for mobile data).
\item \textbf{Provider.} A subscription with an \cle{ISP} (Internet Service Provider) — e.g.\ CAMTEL, MTN Cameroon, Orange Cameroon — which supplies the connection and assigns an IP address.
\item \textbf{Software.} An operating system with the TCP/IP protocol stack configured, plus application software: a web browser (Chrome, Firefox), an email client, etc.
\item \textbf{Account.} Login credentials (username/password or an activated SIM with a data plan) to authenticate on the ISP's network.
\item \textbf{Exam tip.} If the question says ``\emph{four} requirements'', give exactly four, one per category, and name real examples to secure full marks.
\end{enumerate}
\end{methode}

\begin{exoresolu}[Equipping a school computer room]
\textbf{Statement.}\\
GBHS Kumba has just received 20 computers and wants to connect them to the Internet.\\
(a) State four requirements the school must fulfil to obtain an Internet connection. \hfill (4 marks)\\
(b) The school subscribes to an ISP. Define the acronym \emph{ISP} and give two examples of ISPs operating in Cameroon. \hfill (3 marks)\\
(c) Explain the role of the modem in the connection. \hfill (2 marks)
\tcblower
\textbf{Detailed solution.}\\
(a) \textbf{Four requirements.}
\begin{enumerate}
\item \textbf{Hardware:} the computers, each with a network interface (Ethernet/Wi-Fi), plus a modem-router (e.g.\ an ADSL router) and the linking medium (telephone line or SIM-equipped 4G router).
\item \textbf{An ISP subscription:} a contract with an Internet Service Provider that supplies the access and an IP address.
\item \textbf{Software:} an operating system with TCP/IP configured and a web browser installed on each machine.
\item \textbf{An account / authentication:} the username and password (or activated data SIM) used to log on to the ISP's network.
\end{enumerate}

(b) \textbf{ISP} means \textbf{Internet Service Provider}: a company that provides individuals and organisations with access to the Internet, usually for a subscription fee. Examples in Cameroon: \textbf{CAMTEL} (the historic fixed-line operator), \textbf{MTN Cameroon} and \textbf{Orange Cameroon} (mobile Internet), also Ringo or YooMee-type operators. \emph{(1 mark for the definition, 1 mark per example, up to 2.)}

(c) \textbf{Role of the modem.} The modem (\emph{modulator--demodulator}) converts the digital signals of the computer into signals suitable for the transmission line (modulation) when sending, and converts the incoming line signals back into digital data (demodulation) when receiving. It is the indispensable interface between the local equipment and the ISP's line.
\end{exoresolu}

\courssub{Skill 5 — Distinguishing Dial-up, ADSL and Fibre; Understanding Bandwidth}

\begin{definition}[Bandwidth / Bit Rate]
\cle{Bandwidth} (or bit rate) is the maximum amount of data that can be transmitted per unit of time. Units: $\mathrm{bit/s}$, $\mathrm{kbit/s}=10^3\,\mathrm{bit/s}$, $\mathrm{Mbit/s}=10^6\,\mathrm{bit/s}$, $\mathrm{Gbit/s}=10^9\,\mathrm{bit/s}$.
\end{definition}

\begin{propriete}[Generations of Fixed-Line Access]
\begin{itemize}
\item \textbf{Dial-up (PSTN):} analogue modem, up to 56 kbit/s, occupies the voice band $\rightarrow$ blocks phone calls.
\item \textbf{ADSL (Asymmetric Digital Subscriber Line):} digital, uses higher frequencies on the same copper pair, allows simultaneous voice + Internet (splitter/filter), \emph{asymmetric}: downstream $>$ upstream.
\item \textbf{Fibre (FTTH -- Fibre To The Home):} light pulses in glass fibre, very high speed (100 Mbit/s -- 10 Gbit/s+), symmetric options, immune to electromagnetic interference, distance-insensitive.
\end{itemize}
\end{propriete}

\begin{methode}[How to handle questions on access speeds and media]
\begin{enumerate}
\item \textbf{Recall the units.} Bit rate in bit/s: $1\ \mathrm{kbit/s} = 10^3\ \mathrm{bit/s}$, $1\ \mathrm{Mbit/s} = 10^6\ \mathrm{bit/s}$, $1\ \mathrm{Gbit/s} = 10^9\ \mathrm{bit/s}$.
\item \textbf{Know the generations.} Dial-up (analogue, 56 kbit/s, blocks phone) $\rightarrow$ ADSL (digital, same copper, simultaneous phone+Internet, asymmetric) $\rightarrow$ Fibre (light, very high speed, symmetric possible).
\item \textbf{Convert when needed.} Download time $t = \dfrac{\text{file size in bits}}{\text{bit rate in bit/s}}$. Remember $1\ \text{byte} = 8\ \text{bits}$.
\item \textbf{Interpret ``asymmetric''.} In ADSL, the A means the downstream (download) rate exceeds the upstream (upload) rate — suitable for home users who download more than they send.
\end{enumerate}
\end{methode}

\begin{exoresolu}[From dial-up to fibre — bandwidth calculation]
\textbf{Statement.}\\
(a) Explain why ADSL allows a subscriber to make a phone call and browse the Internet at the same time on the same line, unlike a dial-up connection. \hfill (2 marks)\\
(b) What does the ``A'' in ADSL stand for, and what does it mean for the user? \hfill (2 marks)\\
(c) A student downloads a 60 MB video file. Calculate the theoretical download time (i) on a 2 Mbit/s ADSL line and (ii) on a 100 Mbit/s fibre connection. (Take $1\ \mathrm{MB} = 8 \times 10^6$ bits.) \hfill (4 marks)
\tcblower
\textbf{Detailed solution.}\\
(a) \textbf{Simultaneous use.} A dial-up modem transmits data by occupying the same frequency band as the voice, so the line is busy and calls are impossible during a connection. ADSL, on the contrary, splits the copper line into separate frequency bands (using a splitter/filter): one band for voice (0--4 kHz), others for data. Voice and data travel simultaneously without interfering.

(b) \textbf{Meaning of ``A''.} ``A'' stands for \textbf{Asymmetric}. It means the download (downstream) bit rate is higher than the upload (upstream) bit rate. For the user, web pages, videos and downloads arrive fast, but sending large files (e.g.\ uploading a video) is noticeably slower.

(c) \textbf{Download times.}\\
File size in bits: $60\ \mathrm{MB} \times 8 \times 10^6 = 4.8 \times 10^8$ bits.
\begin{align*}
\text{(i) ADSL: } t &= \frac{4.8 \times 10^8\ \text{bits}}{2 \times 10^6\ \text{bit/s}} = 240\ \text{s} = 4\ \text{minutes}.\\
\text{(ii) Fibre: } t &= \frac{4.8 \times 10^8\ \text{bits}}{100 \times 10^6\ \text{bit/s}} = 4.8\ \text{s}.
\end{align*}
\textbf{Comment.} The fibre link is 50 times faster here, which matches the ratio of the bit rates ($100/2 = 50$). In practice, times are slightly longer because of protocol overheads and congestion.
\end{exoresolu}

\courssub{Skill 6 — Analysing Mobile and Satellite Connectivity (Cameroonian Context)}

\begin{definition}[Mobile Network Generations]
\begin{itemize}
\item \textbf{2G (GSM):} voice + SMS, slow data (GPRS/EDGE $\approx$ 100--200 kbit/s).
\item \textbf{3G (UMTS/HSPA):} mobile Internet, few Mbit/s.
\item \textbf{4G / LTE:} fast mobile broadband, tens of Mbit/s (100+ Mbit/s peak).
\item \textbf{5G (NR):} very high speed (Gbit/s), low latency ($<$10 ms), massive IoT.
\end{itemize}
\end{definition}

\begin{definition}[Satellite Internet (VSAT)]
A \cle{VSAT} (Very Small Aperture Terminal) dish communicates with a geostationary satellite (GEO, $\approx$ 36\,000 km altitude) or a low-earth-orbit (LEO) constellation. The satellite relays traffic to a ground station (teleport) connected to the Internet backbone.
\end{definition}

\begin{aretenir}[Mobile vs Satellite — Key Trade-offs]
\begin{center}
\begin{tabularx}{\linewidth}{|l|X|X|}
\hline
\rowcolor{popblueL} \textbf{Criterion} & \textbf{Mobile (3G/4G/5G)} & \textbf{Satellite (VSAT)} \\
\hline
Coverage & Good in towns, weak/absent in remote rural areas & Everywhere with sky view (including remote villages, sea) \\
\hline
Equipment cost & Low (modem/router + SIM, few thousand FCFA) & High (dish, transceiver, modem, installation) \\
\hline
Latency & Low (10--50 ms) & High for GEO (500--700 ms), lower for LEO (20--40 ms) \\
\hline
Weather sensitivity & Low & Rain fade (heavy rain degrades signal, especially Ku/Ka band) \\
\hline
Typical use & General public, smartphones, 4G routers & Backup links, remote sites, maritime, emergency \\
\hline
\end{tabularx}
\end{center}
\end{aretenir}

\begin{methode}[How to discuss wireless connectivity in an essay]
\begin{enumerate}
\item \textbf{Mobile networks.} Describe the principle: the phone/router communicates by radio with the nearest base station (cell), which is linked to the operator's core network and then to the Internet. Generations: 2G (voice + SMS, slow data), 3G (mobile Internet), 4G/LTE (fast mobile broadband), 5G (very high speed, low latency).
\item \textbf{Satellite.} A dish (VSAT) communicates with a satellite that relays to a ground station connected to the Internet. Advantage: coverage everywhere, even remote villages. Limitations: high cost of equipment, latency (signal travels to orbit and back), sensitivity to heavy rain.
\item \textbf{Structure an essay.} Introduction (define connectivity) $\rightarrow$ body (one paragraph per technology with advantages + limitations) $\rightarrow$ conclusion (match to need, e.g.\ bridging the digital divide in rural Cameroon).
\item \textbf{Use local colour accurately.} MTN and Orange provide 3G/4G in most Cameroonian towns; CAMTEL operates the fibre backbone; ANTIC regulates and promotes secure use of ICTs.
\end{enumerate}
\end{methode}

\begin{exoresolu}[Connecting a rural health centre]
\textbf{Statement.}\\
A health centre in a village of the Far North Region has no telephone line, no fibre, and weak but usable 3G coverage. The Ministry of Health wants to connect it to the Internet so that patient records can be sent to the regional hospital.\\
(a) Name two technologies that could provide the connection. \hfill (2 marks)\\
(b) For each technology, give one advantage and one limitation in this context. \hfill (4 marks)\\
(c) Which technology do you recommend if the budget is very limited? Justify. \hfill (2 marks)
\tcblower
\textbf{Detailed solution.}\\
(a) \textbf{Two possible technologies.} (i) The \textbf{mobile network (3G/4G)} via a data SIM and a 3G/4G router or modem; (ii) a \textbf{satellite (VSAT)} connection with a dish.

(b) \textbf{Advantages and limitations.}
\begin{center}
\begin{tabularx}{\linewidth}{|l|X|X|}
\hline
\rowcolor{popblueL} \textbf{Technology} & \textbf{Advantage} & \textbf{Limitation} \\
\hline
Mobile 3G/4G & Cheap equipment (a simple modem/router and SIM); quick to install & Weak coverage means slow, unstable speeds; data bundles costly for heavy use \\
\hline
Satellite (VSAT) & Works everywhere, independent of local infrastructure; stable bandwidth & Expensive dish and subscription; noticeable latency (GEO); performance can degrade in heavy rain \\
\hline
\end{tabularx}
\end{center}

(c) \textbf{Recommendation.} With a very limited budget, the \textbf{3G/4G mobile connection} is recommended: the equipment costs only a few thousand FCFA (a modem and SIM from MTN or Orange), installation is immediate, and the weak signal can be improved with an external directional antenna. Satellite would give better quality but its equipment and subscription costs are far beyond a small budget.
\end{exoresolu}

\coursec{Technical Requirements, Addressing and Security for Internet Access}

\courssub{Skill 7 — Drawing and Interpreting a Connection Diagram}

\begin{methode}[How to sketch the path from a home computer to a web server]
\begin{enumerate}
\item \textbf{Place the elements in order.} Computer $\rightarrow$ (router/switch) $\rightarrow$ modem $\rightarrow$ transmission medium (copper line, fibre, radio) $\rightarrow$ ISP's network $\rightarrow$ Internet backbone $\rightarrow$ destination server.
\item \textbf{Label each link} with the technology used (Ethernet/Wi-Fi at home, ADSL/4G/fibre on the access link, fibre on the backbone).
\item \textbf{Mark the roles.} The modem adapts the signal; the router directs packets; the ISP assigns the IP address and routes traffic; DNS servers translate names (e.g.\ \emph{www.minesec.gov.cm}) into IP addresses.
\item \textbf{In an exam}, a clear, labelled diagram with these five or six elements earns full marks; arrows must show the direction of the data flow.
\end{enumerate}
\end{methode}

\begin{exoresolu}[Diagram of a home connection]
\textbf{Statement.}\\
Draw a labelled diagram showing how a student's laptop in Yaoundé, connected to a 4G router, reaches a web server hosted in Europe. Identify four elements of the diagram and state the role of each. \hfill (6 marks)
\tcblower
\textbf{Detailed solution.}\\

\begin{popfigure}
\begin{tikzpicture}[node distance=1.1cm and 1.5cm,
box/.style={draw=popblue, fill=popblueL, rounded corners, minimum height=0.9cm, align=center, font=\small, text width=2.1cm}]
\node[box] (pc) {Laptop (Wi-Fi)};
\node[box, right=of pc] (router) {4G router / modem};
\node[box, right=of router] (bts) {Base station (4G cell)};
\node[box, right=of bts] (isp) {ISP core network};
\node[box, right=of isp] (backbone) {Internet backbone (fibre)};
\node[box, right=of backbone] (server) {Web server (Europe)};
\draw[-latex, thick, popdark] (pc) -- (router);
\draw[-latex, thick, popdark] (router) -- node[above, font=\scriptsize]{radio} (bts);
\draw[-latex, thick, popdark] (bts) -- (isp);
\draw[-latex, thick, popdark] (isp) -- (backbone);
\draw[-latex, thick, popdark] (backbone) -- (server);
\end{tikzpicture}
\legende{Path of a request from a laptop in Yaoundé to a web server in Europe.}
\end{popfigure}

\textbf{Roles of four elements (any four).}
\begin{enumerate}
\item \textbf{4G router/modem:} adapts the laptop's data to the radio link (modulation) and shares the connection with local devices (Wi-Fi AP + router).
\item \textbf{Base station (BTS):} receives the radio signal from the router and forwards the traffic into the operator's core network via fibre/ microwave backhaul.
\item \textbf{ISP core network:} authenticates the subscriber (AAA), assigns an IP address (dynamic or static), routes packets towards the Internet; DNS servers translate the website name into an IP address.
\item \textbf{Internet backbone:} high-capacity fibre links (including submarine cables landing on the Cameroonian coast: WACS, SAT-3, ACE) that carry the traffic across continents to the destination server.
\end{enumerate}
\emph{Marking guide: 2 marks for a correct, ordered, labelled diagram; 1 mark per correctly stated role (maximum 4).}
\end{exoresolu}

\begin{experience}[Guided Activity: Trace Your Own Connection]
\textbf{Objective:} Understand the path your data takes.
\begin{enumerate}
\item Open a terminal (Command Prompt on Windows, Terminal on Linux/macOS).
\item Run \texttt{tracert www.google.com} (Windows) or \texttt{traceroute www.google.com} (Linux/macOS).
\item Observe the list of hops (routers). The first hop is your router, the next few belong to your ISP (CAMTEL, MTN, Orange), then you see backbone routers, possibly international ones.
\item Note the round-trip times (RTT) in milliseconds. A jump in RTT often indicates a long-distance link (e.g., submarine cable).
\end{enumerate}
\textbf{Reflection:} Why do some hops show ``* * *''? (Answer: routers configured not to respond to ICMP TTL-exceeded messages for security.)
\end{experience}

\begin{attention}[Security Basics for Internet Access]
\begin{itemize}
\item \textbf{Change default router passwords.} Many attacks exploit unchanged admin credentials on home routers.
\item \textbf{Use WPA2/WPA3 encryption} on Wi-Fi, never WEP or open networks.
\item \textbf{Keep firmware updated.} ISPs (CAMTEL, MTN, Orange) push updates to CPE (Customer Premises Equipment); do not block them.
\item \textbf{Beware of public Wi-Fi.} Use a VPN when accessing sensitive services (banking, school portal) on open hotspots.
\item \textbf{ANTIC's role.} In Cameroon, ANTIC operates the \cle{CERT} (Computer Emergency Response Team) and promotes cybersecurity awareness; report incidents to \texttt{cert@antic.cm}.
\end{itemize}
\end{attention}

\begin{exercice}[Chapter Review Questions]
\begin{enumerate}
\item \textbf{History.} (a) State the year ARPANET adopted TCP/IP. (b) Who invented the World Wide Web and where? (c) Explain why the Internet and the Web are not the same thing.
\item \textbf{Packet switching.} A 10 MB file is sent over a packet-switched network. (a) What happens to the file at the sender's TCP layer? (b) If packets arrive out of order, how does TCP handle it?
\item \textbf{Connectivity.} A small business in Douala needs a stable 50 Mbit/s symmetric connection for video conferencing and cloud backups. The building has a telephone line and is covered by 4G and fibre. Recommend the best technology and justify with three technical arguments.
\item \textbf{Requirements.} List the four categories of requirements for a standalone PC to access the Internet. For each, give a concrete example (hardware model, ISP name, software name, credential type).
\item \textbf{Bandwidth calculation.} A 250 MB file is downloaded on a 20 Mbit/s link. Calculate the theoretical minimum download time in seconds. (Use $1\ \mathrm{MB}=8\times10^6$ bits.)
\item \textbf{Mobile vs Satellite.} Compare 4G and GEO satellite for a remote school in the East Region. Use a table with four criteria.
\item \textbf{Diagram.} Draw the connection path for a smartphone in Bafoussam using mobile data to access a server in South Africa. Label five elements and state the role of the base station and the ISP core network.
\end{enumerate}
\end{exercice}

\begin{corrige}[Chapter Review Questions]
\begin{enumerate}
\item (a) 1983 (1 January). (b) Tim Berners-Lee at CERN, Geneva (1989--1991). (c) Internet = infrastructure (network of networks, TCP/IP); Web = service (HTTP/HTML/URLs) running on the Internet.
\item (a) TCP segments the file into numbered packets, adds headers with destination IP. (b) TCP buffers packets, reorders them by sequence number, requests retransmission of missing ones, then delivers the reassembled file to the application.
\item \textbf{Recommendation: Fibre (FTTH).} Arguments: (1) Symmetric 50+ Mbit/s easily available; (2) Dedicated, stable bandwidth not shared with neighbours; (3) Low latency ($<$5 ms) essential for video conferencing; (4) Unlimited data plans common, unlike 4G bundles.
\item (1) Hardware: PC + Ethernet card + fibre ONT (e.g., Huawei HG8245). (2) ISP: CAMTEL fibre subscription. (3) Software: Windows 11 (TCP/IP stack) + Firefox browser. (4) Authentication: PPPoE username/password provided by CAMTEL.
\item File = $250 \times 8 \times 10^6 = 2 \times 10^9$ bits. Rate = $20 \times 10^6$ bit/s. $t = 2 \times 10^9 / 20 \times 10^6 = 100$ seconds.
\item Table: Coverage (4G: town only; Satellite: everywhere), Cost (4G: low; Satellite: high), Latency (4G: low; Satellite: high $\approx$600 ms), Weather (4G: robust; Satellite: rain fade).
\item Diagram: Smartphone $\rightarrow$ (radio) $\rightarrow$ BTS $\rightarrow$ (fibre/microwave) $\rightarrow$ ISP Core $\rightarrow$ (backbone/submarine cable) $\rightarrow$ Server (ZA). Roles: BTS = radio access + backhaul; ISP Core = authentication, IP assignment, routing, DNS.
\end{enumerate}
\end{corrige}

\newpage

\coursec{Exercises}

\courssub{I check my knowledge}

\begin{exercice}[Multiple-choice questions]
For each question, choose the single correct answer.
\begin{enumerate}
\item The protocol suite that replaced NCP on ARPANET on 1 January 1983 is:
\begin{enumerate}
\item HTTP/HTTPS \item TCP/IP \item FTP/SMTP \item OSI/RM
\end{enumerate}
\item The World Wide Web was invented in 1989--1991 by:
\begin{enumerate}
\item Vint Cerf \item Tim Berners-Lee \item Leonard Kleinrock \item Paul Baran
\end{enumerate}
\item Which technology typically offers the \emph{lowest} latency for a home connection in Yaound\'e?
\begin{enumerate}
\item Satellite (GEO) \item Dial-up \item Fibre to the home (FTTH) \item 2G mobile data
\end{enumerate}
\item NAT (Network Address Translation) is normally performed by:
\begin{enumerate}
\item the switch \item the router \item the repeater \item the web browser
\end{enumerate}
\end{enumerate}
\end{exercice}

\begin{exercice}[True or false — justify]
State whether each statement is true or false, and justify your answer in one or two sentences.
\begin{enumerate}
\item ARPANET used circuit switching to guarantee a dedicated path for every message.
\item A private IP address such as 192.168.1.10 can be routed directly on the public Internet.
\item Doubling bandwidth always halves the time needed to load a web page.
\item The landing of submarine cables such as ACE and 2Africa on the Cameroonian coast can improve both redundancy and affordability of Internet access.
\end{enumerate}
\end{exercice}

\begin{exercice}[Definitions]
Define each of the following terms in one or two clear sentences, giving an example where appropriate:
\begin{enumerate}
\item packet switching;
\item bandwidth;
\item latency;
\item ISP (Internet Service Provider);
\item IP address.
\end{enumerate}
\end{exercice}

\begin{exercice}[Direct calculations]
\begin{enumerate}
\item A file of 600 MB is downloaded over a link with an effective throughput of 40 Mbps. Calculate the download time in seconds, then in minutes (recall: 1 byte = 8 bits).
\item A video-conference application requires a round-trip latency below 150 ms. A GEO satellite link has a round-trip latency of about 600 ms. State, with justification, whether this link is suitable.
\end{enumerate}
\end{exercice}

\courssub{I practise}

\begin{exercice}[Sizing a school connection]
A secondary school in Bafoussam plans to stream 5 concurrent 720p video lessons, each requiring 3 Mbps.
\begin{enumerate}
\item Calculate the total bandwidth needed for the streams.
\item Add 20\,\% overhead for protocol traffic and other uses, and determine the minimum symmetric fibre plan (in Mbps) the school should subscribe to.
\item Explain why a \emph{symmetric} plan is recommended if teachers must also upload recorded lessons.
\end{enumerate}
\end{exercice}

\begin{exercice}[Labelling a home network]
Draw (or copy from your course) the schematic of a typical home network: the ISP line enters the house and serves a desktop PC by cable, plus a laptop and two smartphones wirelessly.
\begin{enumerate}
\item Label the following devices: modem, router, switch, Wi-Fi access point.
\item Indicate clearly where NAT occurs and justify your answer.
\item State one advantage of connecting the desktop PC by cable rather than Wi-Fi.
\end{enumerate}
\end{exercice}

\begin{exercice}[Comparing connectivity technologies]
Copy and complete the following comparison table for a cybercaf\'e owner in Bamenda choosing an Internet connection. Use the terms \emph{low / medium / high} for cost and latency, and give a typical bandwidth range.

\begin{center}
\begin{tabularx}{\linewidth}{|l|X|X|X|}\hline
\rowcolor{popblueL} \textbf{Technology} & \textbf{Typical bandwidth} & \textbf{Latency} & \textbf{Installation cost} \\ \hline
ADSL & & & \\ \hline
4G mobile & & & \\ \hline
FTTH (fibre) & & & \\ \hline
VSAT satellite & & & \\ \hline
\end{tabularx}
\end{center}

Then recommend one technology for the cybercaf\'e, justifying your choice with two arguments.
\end{exercice}

\begin{exercice}[Bandwidth versus latency]
\begin{enumerate}
\item Explain, in your own words, the difference between bandwidth and latency, using the analogy of a road between Douala and Yaound\'e.
\item Give one real-world example of an application where a connection with \emph{high bandwidth but high latency} still degrades performance, and explain why.
\item Name one application for which high latency matters little but high bandwidth is essential.
\end{enumerate}
\end{exercice}

\courssub{I prepare for the GCE}

\begin{exercice}[Essay: submarine cables and Cameroon \hfill (20 marks)]
``The landing of submarine fibre-optic cables such as ACE and 2Africa on the Cameroonian coast transforms the national Internet landscape.''Discuss this statement. Your essay should:
\begin{enumerate}
\item briefly recall the historical development of the Internet from ARPANET to the commercial Web; \textbf{(4 marks)}
\item explain the role of submarine cables in international connectivity and in redundancy (what happens when a country depends on a single cable?); \textbf{(6 marks)}
\item analyse the expected effects on affordability, quality of service and the digital economy in Cameroon (e-government, mobile money, online education); \textbf{(6 marks)}
\item identify two remaining obstacles that prevent all Cameroonians from benefiting fully, and propose a solution for each. \textbf{(4 marks)}
\end{enumerate}
\end{exercice}

\begin{exercice}[Case study: a business in Douala \hfill (15 marks)]
A small import-export business in Bonab\'eri, Douala, uses an ADSL line to connect its 8 employees. The accountant complains that the VPN connection to the bank's server drops several times per day, especially when it rains, and that large file uploads are very slow.
\begin{enumerate}
\item State three structured troubleshooting steps the technician should perform first, and explain what each step can reveal. \textbf{(6 marks)}
\item Give two technical reasons why ADSL can suffer from instability and low upload speed. \textbf{(4 marks)}
\item Propose a long-term upgrade path (technology, equipment, and a redundancy solution), justifying each choice. \textbf{(5 marks)}
\end{enumerate}
\end{exercice}

\begin{exercice}[Structured problem: connecting a new school \hfill (20 marks)]
A new Government High School in a rural divisional headquarters must be connected to the Internet for its computer laboratory of 40 machines. The town has 4G mobile coverage but no fibre; the national electricity supply is unstable.
\begin{enumerate}
\item List the hardware and software requirements needed to connect the laboratory and share the connection among the 40 machines. \textbf{(5 marks)}
\item Compare two candidate access technologies (4G router vs.\ VSAT satellite) in terms of bandwidth, latency, cost and reliability, and make a justified recommendation. \textbf{(6 marks)}
\item Explain how IP addressing should be organised inside the laboratory (private addresses, DHCP, NAT). \textbf{(5 marks)}
\item Propose two security measures to protect students and the school's data, and one measure to cope with power cuts. \textbf{(4 marks)}
\end{enumerate}
\end{exercice}

\newpage

\coursec{Solutions to the exercises}

\begin{corrige}[Exercise 1 — Multiple-choice questions]
\begin{enumerate}
\item \textbf{Answer: b) TCP/IP.} On 1 January 1983, ARPANET switched from the NCP protocol to the TCP/IP suite, an event often called the ``flag day''. This date is considered the birth of the modern Internet. HTTP (a) is a Web application protocol (1991), FTP/SMTP (c) are application protocols that run \emph{on top of} TCP/IP, and OSI/RM (d) is a theoretical reference model, never deployed on ARPANET.
\item \textbf{Answer: b) Tim Berners-Lee.} Working at CERN, he proposed the Web in 1989 and implemented HTML, HTTP and the first browser/server around 1990--1991. Vint Cerf co-designed TCP/IP, Kleinrock developed the theory of packet switching, and Baran invented the concept of distributed packet-switched networks.
\item \textbf{Answer: c) Fibre to the home (FTTH).} Light propagates in fibre with very low delay (typically a few ms to the ISP). GEO satellite (a) imposes about 500--600 ms round trip, dial-up (b) has high latency and tiny bandwidth, and 2G (d) is slow with high latency.
\item \textbf{Answer: b) the router.} NAT translates private addresses (e.g.\ 192.168.1.x) into the single public address assigned by the ISP; this is a Layer-3 function performed by the router (often integrated in the ISP box). A switch works at Layer 2, a repeater at Layer 1, and the browser is application software.
\end{enumerate}
\end{corrige}

\begin{corrige}[Exercise 2 — True or false]
\begin{enumerate}
\item \textbf{False.} ARPANET was precisely the first large network to use \cle{packet switching}: messages are cut into packets routed independently, with no dedicated path. Circuit switching (like the old telephone network) was exactly what Baran and Kleinrock wanted to avoid.
\item \textbf{False.} Addresses in the private ranges (10.0.0.0/8, 172.16.0.0/12, 192.168.0.0/16) are not routable on the public Internet; Internet routers drop them. They reach the Internet only through NAT on the router, which substitutes a public address.
\item \textbf{False.} Doubling bandwidth halves the \emph{transfer time of the data volume}, but the total loading time also includes latency, DNS resolution, server response time and processing, which do not depend on bandwidth. Hence the gain is real but never exactly a factor of two.
\item \textbf{True.} Each additional submarine cable provides an alternative international route: if one cable is cut, traffic can be rerouted (redundancy), and increased competition/capacity between cable systems tends to lower wholesale bandwidth prices, which can translate into cheaper retail access.
\end{enumerate}
\end{corrige}

\begin{corrige}[Exercise 3 — Definitions]
\begin{enumerate}
\item \textbf{Packet switching:} a data-transmission technique in which a message is divided into small blocks called \cle{packets}, each carrying addressing information and routed independently through the network, then reassembled at the destination. Example: an e-mail sent from Douala to Bamenda is split into packets that may follow different paths.
\item \textbf{Bandwidth:} the maximum quantity of data a link can carry per unit of time, measured in bits per second (bps, Mbps, Gbps). Example: a 20 Mbps FTTH subscription.
\item \textbf{Latency:} the time delay for a data unit to travel from source to destination (often measured as round-trip time), expressed in milliseconds. Example: about 600 ms round trip on a GEO satellite link.
\item \textbf{ISP (Internet Service Provider):} a company that sells Internet access to customers and routes their traffic to the global Internet. Examples in Cameroon: MTN Cameroon, Orange Cameroon, CAMTEL, Nexttel.
\item \textbf{IP address:} a unique numerical identifier assigned to a device on an IP network, used to address and route packets to it. Example: 192.168.1.10 (IPv4 private address).
\end{enumerate}
\end{corrige}

\begin{corrige}[Exercise 4 — Direct calculations]
\begin{enumerate}
\item Convert the file size to bits:
\[ 600\ \text{MB} = 600 \times 8 = 4800\ \text{Mb (megabits)}. \]
The download time is
\[ t = \dfrac{4800\ \text{Mb}}{40\ \text{Mbps}} = 120\ \text{s} = \dfrac{120}{60} = 2\ \text{min}. \]
\item \textbf{The GEO link is not suitable.} Its round-trip latency (about 600 ms) is four times the maximum tolerated (150 ms): $600 > 150$. Even with abundant bandwidth, the delay would cause awkward pauses and echoes in the conversation, making interactive video-conferencing unusable. A terrestrial link (fibre, 4G) with latency below 150 ms must be used.
\end{enumerate}
\end{corrige}

\begin{corrige}[Exercise 5 — Sizing a school connection]
\begin{enumerate}
\item Total bandwidth for the streams:
\[ 5 \times 3\ \text{Mbps} = 15\ \text{Mbps}. \]
\item Adding 20\,\% overhead:
\[ 15 \times 1.20 = 18\ \text{Mbps}. \]
The school should therefore subscribe to a symmetric fibre plan of \textbf{at least 18 Mbps}; in practice the next commercial tier, e.g.\ \textbf{20 Mbps symmetric}, would be chosen.
\item Streaming only consumes \emph{downstream} bandwidth, but uploading recorded lessons consumes \emph{upstream} bandwidth. An asymmetric plan (e.g.\ ADSL) typically offers an upload rate 5 to 10 times lower than the download rate, so publishing videos would be painfully slow and could saturate the uplink, degrading the whole connection (acknowledgements of downloads also use the uplink). A \emph{symmetric} plan guarantees the same capacity in both directions, so teachers can upload lessons while students stream without congestion.
\end{enumerate}
\end{corrige}

\begin{corrige}[Exercise 6 — Labelling a home network]
\begin{enumerate}
\item Schematic of the typical home network:

\begin{popfigure}
\begin{center}
\begin{tikzpicture}[scale=0.95, every node/.style={font=\small}]
\node[draw=popdark, rounded corners, fill=popblueL, minimum width=2.2cm, minimum height=0.9cm] (isp) at (0,0) {ISP line};
\node[draw=popdark, rounded corners, fill=popgreenL, minimum width=2.2cm, minimum height=0.9cm] (modem) at (3.2,0) {Modem};
\node[draw=popdark, rounded corners, fill=popblueL, minimum width=2.6cm, minimum height=0.9cm] (router) at (6.6,0) {Router (NAT)};
\node[draw=popdark, rounded corners, fill=popgreenL, minimum width=2.2cm, minimum height=0.9cm] (switch) at (10,1.6) {Switch};
\node[draw=popdark, rounded corners, fill=popgreenL, minimum width=2.6cm, minimum height=0.9cm] (ap) at (10,-1.6) {Wi-Fi access point};
\node[draw=popdark, rounded corners, fill=white, minimum width=1.9cm, minimum height=0.8cm] (pc) at (13.4,1.6) {Desktop PC};
\node[draw=popdark, rounded corners, fill=white, minimum width=1.7cm, minimum height=0.8cm] (lap) at (13.4,-0.9) {Laptop};
\node[draw=popdark, rounded corners, fill=white, minimum width=1.9cm, minimum height=0.8cm] (ph) at (13.4,-2.4) {2 smartphones};
\draw[thick, popdark] (isp) -- (modem);
\draw[thick, popdark] (modem) -- (router);
\draw[thick, popdark] (router.east) -- (8.6,0) |- (switch.west);
\draw[thick, popdark] (router.east) -- (8.6,0) |- (ap.west);
\draw[thick, popdark] (switch) -- (pc);
\draw[thick, poporange, dashed] (ap) -- (lap);
\draw[thick, poporange, dashed] (ap) -- (ph);
\end{tikzpicture}
\end{center}
\legende{Home network: ISP line $\rightarrow$ modem $\rightarrow$ router $\rightarrow$ switch (wired PC) and Wi-Fi access point (laptop, smartphones).}
\end{popfigure}

In most real installations the modem, router, switch and access point are combined in a single box supplied by the ISP.
\item \textbf{NAT occurs in the router.} The router is the boundary between the private LAN (addresses such as 192.168.1.x) and the public Internet: it translates each private source address into its own public address (with port numbers to distinguish the devices) before forwarding packets to the ISP, and performs the reverse translation for replies.
\item Advantage of the cable for the desktop PC: a \textbf{more stable and faster connection}, immune to radio interference, walls and distance, with lower latency than Wi-Fi (also slightly more secure against eavesdropping).
\end{enumerate}
\end{corrige}

\begin{corrige}[Exercise 7 — Comparing connectivity technologies]
Completed table (typical values, to be adapted to local offers):

\begin{center}
\begin{tabularx}{\linewidth}{|l|X|X|X|}\hline
\rowcolor{popblueL} \textbf{Technology} & \textbf{Typical bandwidth} & \textbf{Latency} & \textbf{Installation cost} \\ \hline
ADSL & 1--20 Mbps & low to medium & low \\ \hline
4G mobile & 5--50 Mbps & medium & low \\ \hline
FTTH (fibre) & 20 Mbps--1 Gbps & low & medium to high \\ \hline
VSAT satellite & 2--30 Mbps & high (500--700 ms) & high \\ \hline
\end{tabularx}
\end{center}

\textbf{Recommendation: FTTH (fibre)} if it is available at the cybercaf\'e's location in Bamenda, for two reasons:
\begin{enumerate}
\item \textbf{Quality of service:} fibre offers the highest bandwidth and the lowest latency, so many customers can browse, stream and make video calls simultaneously without congestion — essential for a business whose revenue depends on connection quality.
\item \textbf{Reliability and cost per Mbps:} once installed, fibre is stable (unaffected by weather, unlike satellite or radio) and its cost per Mbps is the lowest, which improves the profitability of the cybercaf\'e.
\end{enumerate}
(If fibre is not yet available in the neighbourhood, a 4G router with a generous data plan is the acceptable fallback; VSAT is the last resort because of its high latency and cost.)
\end{corrige}

\begin{corrige}[Exercise 8 — Bandwidth versus latency]
\begin{enumerate}
\item Imagine the Douala--Yaound\'e road. \textbf{Bandwidth} is the \emph{number of lanes}: the more lanes, the more vehicles (data) can travel at the same time. \textbf{Latency} is the \emph{time a single vehicle needs to go from Douala to Yaound\'e}: even a six-lane motorway does not reduce the travel time of one car. Similarly, a link can carry huge volumes per second (high bandwidth) while each individual packet still takes a long time to arrive (high latency).
\item Example: an \textbf{online video-conference or VoIP call over a GEO satellite link} with high bandwidth but high latency ($\approx 600$ ms round trip). The link can carry the video stream easily, but every reply arrives with a noticeable delay, so speakers constantly talk over each other: interactivity is destroyed because latency, not bandwidth, governs the responsiveness of the conversation.
\item \textbf{Downloading or streaming a large video file} (e.g.\ watching a recorded film): the user only cares that the whole file arrives quickly, so high bandwidth is essential, while a latency of a few hundred milliseconds at the start is barely noticeable.
\end{enumerate}
\end{corrige}

\begin{corrige}[Exercise 9 — Essay: submarine cables and Cameroon (20 marks)]
\textbf{Indicative mark scheme and expected content.}

\textbf{(a) Historical development (4 marks).} Award 1 mark per correct idea, up to 4:
\begin{itemize}
\item 1960s: packet switching theorised (Baran, Kleinrock); 1969: ARPANET connects its first nodes (universities/research in the USA).
\item 1974: TCP/IP designed by Cerf and Kahn; 1 January 1983: ARPANET adopts TCP/IP — birth of the Internet; interconnection of networks (NSFNET, etc.).
\item 1989--1991: Tim Berners-Lee invents the Web (HTML, HTTP, URL) at CERN; 1993: Mosaic browser popularises it.
\item 1990s: commercialisation, ISPs, browsers, search engines; 2000s--today: broadband, mobile Internet, social networks, cloud.
\end{itemize}

\textbf{(b) Role of submarine cables (6 marks).} Expected: submarine fibre-optic cables carry the overwhelming majority of intercontinental Internet traffic (1--2 marks); they land at coastal stations (e.g.\ near Douala/Limb\'e/Kribi) and connect the national backbone to the world (1 mark); each additional cable adds \textbf{capacity} and \textbf{redundancy} (1 mark). A country depending on a \emph{single} cable suffers total or severe international disconnection when that cable is cut (ship anchors, fishing, earthquakes), as has happened to several African countries (2 marks); multiple cables allow automatic rerouting and continuity of service (1 mark).

\textbf{(c) Expected effects for Cameroon (6 marks).} Any three well-developed effects, 2 marks each:
\begin{itemize}
\item \textbf{Affordability:} more capacity and competition lower wholesale prices, which can reduce retail subscription and data costs.
\item \textbf{Quality of service:} higher speeds, lower congestion, better streaming and video-conferencing.
\item \textbf{Digital economy:} growth of e-government services, mobile money (MTN MoMo, Orange Money), online education, e-commerce, outsourcing and start-ups; attraction of data centres and investment.
\end{itemize}

\textbf{(d) Two remaining obstacles + solutions (4 marks).} 1 mark per obstacle and 1 per coherent solution, e.g.:
\begin{itemize}
\item \textbf{Last-mile gap:} cables land on the coast but fibre does not reach rural areas $\rightarrow$ invest in the national fibre backbone, mobile broadband (4G/5G) and community access points.
\item \textbf{Cost of devices and data / digital illiteracy:} many households cannot afford smartphones or lack skills $\rightarrow$ subsidise devices, reduce taxes on ICT equipment, digital literacy programmes in schools.
\item (Also acceptable: unstable electricity supply $\rightarrow$ solar-powered infrastructure; weak competition $\rightarrow$ regulatory enforcement.)
\end{itemize}
\textbf{Examiner's expectations:} structured essay (introduction, development, conclusion), correct technical vocabulary, Cameroonian examples, and balanced discussion (benefits \emph{and} limits), not a mere list.
\end{corrige}

\begin{corrige}[Exercise 10 — Case study: a business in Douala (15 marks)]
\textbf{(a) Three structured troubleshooting steps (6 marks — 2 each: step + what it reveals).}
\begin{enumerate}
\item \textbf{Check the physical layer:} inspect cables, connectors, the ADSL splitter/filter and the modem's synchronisation LED; test with another modem if possible. \emph{Reveals:} corroded or loose connectors, water ingress in outdoor joints (consistent with drops when it rains), or a faulty modem.
\item \textbf{Test the line quality:} read the modem's diagnostic values (SNR — signal-to-noise ratio, line attenuation, error counts) and ask the ISP for a line test. \emph{Reveals:} a degraded copper pair, excessive distance from the exchange, or interference — typical causes of instability.
\item \textbf{Isolate the internal network:} connect a single PC directly to the modem, run ping/traceroute to the bank's server and monitor packet loss over time. \emph{Reveals:} whether the problem is inside the company (switch, cabling, overloaded Wi-Fi, misconfigured VPN client) or on the ISP side.
\end{enumerate}

\textbf{(b) Two technical reasons for ADSL instability and low upload speed (4 marks — 2 each).}
\begin{itemize}
\item \textbf{Copper sensitivity:} ADSL uses the old telephone copper pair, whose signal attenuates with distance from the exchange and is vulnerable to moisture, corrosion and electromagnetic interference — hence synchronisation losses, especially when it rains.
\item \textbf{Asymmetry by design:} the ``A'' in ADSL means most of the frequency band is allocated to the downstream; the upstream band is narrow, so upload throughput is typically 5--10 times lower than download, making large file uploads very slow.
\end{itemize}

\textbf{(c) Long-term upgrade path (5 marks).}
\begin{itemize}
\item \textbf{Technology:} migrate to \textbf{FTTH fibre} (or a 4G/5G business router if fibre is unavailable in Bonab\'eri) — symmetric bandwidth, low latency, immunity to weather on the access line (2 marks).
\item \textbf{Equipment:} a business-grade router/firewall with VPN support (hardware VPN acceleration), a manageable switch, and a UPS for the modem/router (1 mark).
\item \textbf{Redundancy:} keep a \textbf{4G failover connection} as backup: the router switches automatically to 4G when the fibre fails, guaranteeing continuity of the VPN to the bank (2 marks).
\end{itemize}
\textbf{Examiner's expectations:} methodical reasoning (physical $\rightarrow$ line $\rightarrow$ network), correct use of terms (SNR, attenuation, asymmetry, failover), and realistic proposals for the Douala context.
\end{corrige}

\begin{corrige}[Exercise 11 — Structured problem: connecting a new school (20 marks)]
\textbf{(a) Hardware and software requirements (5 marks — 1 per item, any five).}
\begin{itemize}
\item An Internet access device: 4G router (or VSAT terminal) with an active ISP subscription/SIM.
\item A switch (48 ports or two 24-port switches) and Ethernet cabling to connect the 40 machines.
\item Optionally one or more Wi-Fi access points for mobile devices.
\item A router/firewall function (often integrated in the 4G router) to share the connection (NAT) and filter content.
\item A DHCP service (usually on the router) to distribute IP addresses automatically.
\item Software: operating systems updated, antivirus, a web browser, and optionally a proxy/filtering software; a server is optional.
\item Electrical protection: UPS or inverter (may be credited here or in (d)).
\end{itemize}

\textbf{(b) 4G router vs VSAT (6 marks).}

\begin{center}
\begin{tabularx}{\linewidth}{|l|X|X|}\hline
\rowcolor{popblueL} \textbf{Criterion} & \textbf{4G router} & \textbf{VSAT satellite} \\ \hline
Bandwidth & 5--50 Mbps, shared with other users & 2--30 Mbps, more stable per contract \\ \hline
Latency & medium (30--80 ms), fine for Web/video & high (500--700 ms), poor for interactivity \\ \hline
Cost & low equipment, affordable data plans & expensive dish + installation + subscription \\ \hline
Reliability & depends on coverage/congestion; weather-sensitive & works anywhere; degraded by heavy rain \\ \hline
\end{tabularx}
\end{center}

\textbf{Recommendation (2 marks):} choose the \textbf{4G router}, because the town already has 4G coverage: it is far cheaper to install and operate, offers lower latency suitable for educational browsing and video, and requires no specialised maintenance. VSAT would only be justified if 4G coverage were absent or unusable.

\textbf{(c) IP addressing organisation (5 marks).}
\begin{itemize}
\item Use a \textbf{private address range}, e.g.\ 192.168.1.0/24, which supports up to 254 hosts — more than enough for 40 machines (1 mark).
\item Configure the router as a \textbf{DHCP server} distributing addresses (e.g.\ 192.168.1.100--192.168.1.200), the default gateway (192.168.1.1) and DNS servers automatically, avoiding manual configuration errors (2 marks).
\item The router performs \textbf{NAT}: all 40 private addresses are translated to the single public address provided by the ISP, so one subscription serves the whole laboratory (2 marks).
\end{itemize}

\textbf{(d) Security and power measures (4 marks).}
\begin{itemize}
\item Security (any two, 1 mark each): a \textbf{firewall} with \textbf{content filtering} to block inappropriate sites and protect students; \textbf{antivirus} software on all machines with regular updates; strong \textbf{passwords/WPA2 or WPA3} on the Wi-Fi; user accounts with limited rights for students.
\item Power (2 marks): install a \textbf{UPS/inverter with batteries} (possibly solar-backed) to keep the router, switch and machines running during cuts and to protect equipment from surges.
\end{itemize}
\textbf{Examiner's expectations:} a complete, coherent and cost-aware solution; correct use of the terms DHCP, NAT, private addressing; and justification of each choice in the rural Cameroonian context (coverage, budget, electricity).
\end{corrige}

\newpage

\coursec{Summary sheet}

\begin{popfigure}
\begin{tikzpicture}[mindmap, grow cyclic, every node/.style={concept, text=white, font=\small\bfseries},
root concept/.append style={concept color=popdark, font=\large\bfseries},
level 1/.append style={level distance=3.6cm, sibling angle=72},
level 2/.append style={level distance=2.4cm, sibling angle=45, font=\scriptsize\bfseries}]
\node[root concept]{History of the Internet}
  child[concept color=popblue]{ node{Origins: ARPANET}
    child{ node{1969, 4 nodes} }
    child{ node{Packet switching} }
  }
  child[concept color=poporange]{ node{TCP/IP}
    child{ node{1974 Cerf \& Kahn} }
    child{ node{1 Jan 1983 flag day} }
  }
  child[concept color=popgreen]{ node{Web Revolution}
    child{ node{1989--91 Berners-Lee} }
    child{ node{HTML, URL, HTTP} }
    child{ node{Browsers: Mosaic} }
  }
  child[concept color=poppurple]{ node{Modern Internet}
    child{ node{Web 2.0, mobile} }
    child{ node{Cameroon: CAMTEL, MTN} }
  }
  child[concept color=poppink]{ node{Connectivity}
    child{ node{Dial-up, ADSL} }
    child{ node{Fibre, 4G, satellite} }
  }
  child[concept color=popgold]{ node{Requirements}
    child{ node{ISP, modem, NIC} }
    child{ node{IP address, browser} }
  };
\end{tikzpicture}
\legende{Mind map of Chapter ICT17: from ARPANET to today's connection technologies.}
\end{popfigure}

\begin{bilan}
\textbf{Key definitions.}
\begin{itemize}
\item \cle{Internet}: a global network of interconnected networks communicating through common protocols (TCP/IP); a \emph{network of networks}.
\item \cle{ARPANET}: the ancestor of the Internet, launched in 1969 by ARPA (US Department of Defense), linking four university computers.
\item \cle{Packet switching}: data is split into small addressed packets routed independently and reassembled at destination (concept by Baran and Davies).
\item \cle{TCP/IP}: the protocol suite adopted on 1 January 1983; \cle{IP} addresses and routes packets, \cle{TCP} guarantees ordered, reliable delivery.
\item \cle{World Wide Web}: system of hypertext documents invented by Tim Berners-Lee at CERN (1989--1991), based on \cle{HTML}, \cle{URL} and \cle{HTTP}. The Web is a \emph{service} of the Internet, not the Internet itself.
\item \cle{ISP} (Internet Service Provider): company that sells Internet access (e.g. CAMTEL, MTN Cameroon, Orange, Nexttel).
\item \cle{Modem}: device that modulates/demodulates signals between the user's equipment and the ISP's line.
\item \cle{Bandwidth}: maximum data rate of a link, measured in bits per second (kbit/s, Mbit/s, Gbit/s).
\end{itemize}

\textbf{Key dates and facts to memorise.}
\begin{itemize}
\item 1969: ARPANET connects UCLA, Stanford, UCSB and Utah; first message ``LO'' (crash during ``LOGIN'').
\item 1971: Ray Tomlinson sends the first networked e-mail and introduces the @ symbol.
\item 1974: Vinton Cerf and Bob Kahn publish the TCP design; 1983: ARPANET switches to TCP/IP; DNS follows in 1984.
\item 1989--1991: Berners-Lee creates the Web; 1993: Mosaic, the first popular graphical browser; 1995: commercial Internet boom.
\item 2000s: broadband (ADSL, fibre), Web 2.0, social networks; 2010s: mobile Internet (3G/4G) dominates, including in Cameroon.
\end{itemize}

\textbf{Connectivity technologies (know advantages/limits).}
\begin{center}
\begin{tabularx}{\linewidth}{|l|X|X|}
\hline
\rowcolor{popblueL} \textbf{Technology} & \textbf{Medium} & \textbf{Remark} \\
\hline
Dial-up & Telephone line & Slow ($\approx 56$ kbit/s), obsolete \\
\hline
ADSL & Telephone copper pair & Broadband, asymmetric speeds \\
\hline
Fibre optic & Glass fibre & Very high speed, low latency \\
\hline
Mobile 3G/4G/5G & Radio waves & Dominant access mode in Cameroon \\
\hline
Satellite & Radio via satellite & Covers remote rural areas \\
\hline
Wi-Fi / WiMAX & Local radio & Wireless local/metropolitan access \\
\hline
\end{tabularx}
\end{center}

\textbf{Requirements for an Internet connection (method).}
\begin{enumerate}
\item Subscribe to an \cle{ISP} and obtain an account/SIM.
\item Hardware: computer or smartphone, \cle{NIC} or wireless adapter, \cle{modem}/router.
\item Software: operating system with TCP/IP stack, \cle{browser} (Chrome, Firefox).
\item Configuration: obtain an \cle{IP address} (static or via DHCP), DNS server addresses.
\item Test the connection (open a web page, use ping).
\end{enumerate}

\textbf{Common mistakes to avoid.}
\begin{itemize}
\item Confusing the \cle{Internet} (infrastructure, 1969/1983) with the \cle{Web} (service, 1989).
\item Saying Berners-Lee ``invented the Internet'' --- he invented the Web.
\item Confusing bandwidth (capacity) with actual speed; forgetting units (bit/s, not byte/s).
\item Forgetting that a modem alone is not enough: an ISP subscription and IP configuration are required.
\end{itemize}

\textbf{GCE tips.}
\begin{itemize}
\item Learn the timeline as a chain: ARPANET $\rightarrow$ TCP/IP $\rightarrow$ DNS $\rightarrow$ Web $\rightarrow$ broadband/mobile.
\item In essays, define terms first, then compare technologies in a table (speed, cost, availability in Cameroon).
\item Cite local examples (CAMTEL fibre, MTN/Orange 4G, mobile money over the Internet) to earn application marks.
\end{itemize}
\end{bilan}

\findecours

\end{document}
